% !TEX root = 第四章_一元函数微分学(答案版).tex
\setcounter{chapter}{3}
\pagestyle{fancy}

\chapter{一元函数微分学}


\section{导数}


\subsection{导数的定义及可微性验证}

%例题4.1.1
\begin{example}
    设 $f(x)=|x|$,则 $f(x)$ 在 $x\neq0$ 处可导,且
    \[
        f'(x)=\dfrac{|x|}{x}.
    \]
\end{example}
\exampleblank

\begin{solution}
    当 $x>0$ 时 $f(x)=x$,当 $x<0$ 时 $f(x)=-x$.因此对 $x\neq0$,
    \[
        f'(x)=
        \begin{cases}
            1,  & x>0, \\
            -1, & x<0,
        \end{cases}
        =\dfrac{|x|}{x}.
    \]
\end{solution}

%例题4.1.2
\begin{example}
    证明:$f(x)=|x|^3$ 在 $x=0$ 处的三阶导数 $f'''(0)$ 不存在.
\end{example}
\exampleblank

\begin{solution}
    分别在原点两侧求导,可得
    \[
        f(x)=
        \begin{cases}
            x^3,  & x\geq0, \\
            -x^3, & x<0,
        \end{cases}
        \qquad
        f''(x)=6|x|.
    \]
    特别地,$f''(0)=0$.但
    \[
        \lim\limits_{h\to0^+}\dfrac{f''(h)-f''(0)}{h}=6,
        \qquad
        \lim\limits_{h\to0^-}\dfrac{f''(h)-f''(0)}{h}=-6.
    \]
    左右极限不相等,故 $f'''(0)$ 不存在.
\end{solution}

%例题4.1.3
\begin{example}
    设
    \[
        f(x)=
        \begin{cases}
            x^2\sin\dfrac{\pi}{x}, & x<0, \\
            A,                     & x=0, \\
            ax^2+b,                & x>0,
        \end{cases}
    \]
    其中 $A,a,b$ 为常数,试问 $A,a,b$ 为何值时 $f(x)$ 在 $x=0$ 处可导,为什么?并求 $f'(0)$.
\end{example}
\exampleblank

\begin{solution}
    可导必连续.由左右极限
    \[
        \lim\limits_{x\to0^-}x^2\sin\dfrac{\pi}{x}=0,
        \qquad
        \lim\limits_{x\to0^+}(ax^2+b)=b,
    \]
    可知连续性要求且只要求 $A=b=0$.此时
    \[
        \lim\limits_{x\to0^-}\dfrac{f(x)-f(0)}{x}
        =\lim\limits_{x\to0^-}x\sin\dfrac{\pi}{x}=0,
    \]
    而
    \[
        \lim\limits_{x\to0^+}\dfrac{f(x)-f(0)}{x}
        =\lim\limits_{x\to0^+}ax=0.
    \]
    因而 $a$ 可取任意实数,$A=b=0$,且 $f'(0)=0$.
\end{solution}

%例题4.1.4
\begin{example}
    设函数 $f(x)$ 在闭区间 $\lbrack 0,1\rbrack$ 上四次连续可微,$f(0)=0,f'(0)=0$,证明函数
    \[
        F(x)=
        \begin{cases}
            \dfrac{f(x)}{x^2}, & 0<x\leq1, \\
            \dfrac{f''(0)}{2}, & x=0,
        \end{cases}
    \]
    在 $\lbrack 0,1\rbrack$ 上二次连续可微.
\end{example}
\exampleblank

\begin{solution}
    由 $f(0)=f'(0)=0$ 及 Taylor 公式的积分余项,对 $x\in\lbrack0,1\rbrack$ 有
    \[
        f(x)=x^2\int_0^1(1-t)f''(tx)\,\mathrm{d}t.
    \]
    因此可将 $F$ 统一写成
    \[
        F(x)=\int_0^1(1-t)f''(tx)\,\mathrm{d}t.
    \]
    当 $x=0$ 时右端等于 $\dfrac{f''(0)}{2}$,与题设定义一致.由于 $f\in C^4(\lbrack0,1\rbrack)$,可在积分号下连续求导两次:
    \[
        F'(x)=\int_0^1t(1-t)f'''(tx)\,\mathrm{d}t,
    \]
    \[
        F''(x)=\int_0^1t^2(1-t)f^{(4)}(tx)\,\mathrm{d}t.
    \]
    上述三个积分都是 $x$ 的连续函数,故 $F\in C^2(\lbrack0,1\rbrack)$.
\end{solution}

%例题4.1.5
\begin{example}
    设 $f(x)$ 在 $x_0$ 的某邻域内有定义.
    \begin{enumerate}
        \item 若 $f(x)$ 在 $x_0$ 处可导,试证
              \[
                  \lim\limits_{h\to0}
                  \dfrac{f(x_0+h)-f(x_0-h)}{2h}=f'(x_0);
              \]
        \item 反之,若上式左端极限存在,是否能推出 $f'(x_0)$ 存在?若结论成立,请证明;不成立请给出反例.
    \end{enumerate}
\end{example}
\exampleblank

\begin{solution}
    \begin{enumerate}
        \item 将差商拆成
              \[
                  \dfrac{f(x_0+h)-f(x_0-h)}{2h}
                  =\dfrac{1}{2}\dfrac{f(x_0+h)-f(x_0)}{h}
                  +\dfrac{1}{2}\dfrac{f(x_0)-f(x_0-h)}{h}.
              \]
              当 $h\to0$ 时,两项分别趋于 $\dfrac{f'(x_0)}{2}$,故原极限为 $f'(x_0)$.
        \item 不能.取 $f(x)=|x-x_0|$,则对任意 $h\neq0$,
              \[
                  \dfrac{f(x_0+h)-f(x_0-h)}{2h}=0,
              \]
              故对称差商的极限存在;但 $f$ 在 $x_0$ 处的左右导数分别为 $-1$ 和 $1$,因而不可导.
    \end{enumerate}
\end{solution}

%例题4.1.6
\begin{example}[\markstar]
    证明 Riemann 函数
    \[
        R(x)=
        \begin{cases}
            \dfrac{1}{q}, & x=\dfrac{p}{q},    \\
            0,            & x=0,1\text{ 或无理数},
        \end{cases}
    \]
    处处不可微.
\end{example}
\exampleblank

\begin{solution}
    先设 $x_0=\dfrac{p_0}{q_0}\in(0,1)$ 为既约有理数.取无理数列 $x_n\to x_0$,则
    \[
        \left|\dfrac{R(x_n)-R(x_0)}{x_n-x_0}\right|
        =\dfrac{1}{q_0|x_n-x_0|}\longrightarrow+\infty,
    \]
    故 $R$ 在 $x_0$ 处不可导.

    再设 $x_0\in(0,1)$ 为无理数.由 Dirichlet 逼近定理,存在无穷多个既约分数 $\dfrac{p}{q}$,使
    \[
        0<\left|x_0-\dfrac{p}{q}\right|<\dfrac{1}{q^2}.
    \]
    沿这些有理点趋于 $x_0$ 时,
    \[
        \left|\dfrac{R\left(\frac{p}{q}\right)-R(x_0)}
        {\frac{p}{q}-x_0}\right|
        =\dfrac{\frac{1}{q}}{\left|\frac{p}{q}-x_0\right|}>q,
    \]
    差商不可能收敛.

    最后,在端点 $0$ 处分别沿无理数和 $x=\dfrac{1}{q}$ 趋近,右差商分别为 $0$ 和 $1$;在端点 $1$ 处分别沿无理数和 $x=\dfrac{q-1}{q}$ 趋近,左差商分别为 $0$ 和 $-1$.因而 $R$ 处处不可微.
\end{solution}

\begin{proposition}[\markstar]
    设函数
    \[
        f(x)=
        \begin{cases}
            x^\alpha\sin\dfrac{1}{x}, & x\neq0, \\
            0,                        & x=0,
        \end{cases}
    \]
    则:
    \begin{enumerate}
        \item $\alpha>0$ 时,在点 $x=0$ 连续;
        \item $\alpha>1$ 时,在点 $x=0$ 处可导;
        \item $\alpha>2$ 时,在点 $x=0$ 处连续可导.
    \end{enumerate}
\end{proposition}

%例题4.1.7
\begin{example}
    求 $f(x)=[x]\sin\pi x$ 的单侧导数,并讨论可微性.
\end{example}
\exampleblank

\begin{solution}
    当 $x\in(m,m+1)$,$m\in\mathbb{Z}$ 时,$[x]=m$,故
    \[
        f'(x)=m\pi\cos\pi x.
    \]
    对整数 $n$,$f(n)=0$.当 $h\to0^+$ 时 $[n+h]=n$,当 $h\to0^-$ 时 $[n+h]=n-1$,因而
    \[
        f'_+(n)=n\pi(-1)^n,
        \qquad
        f'_-(n)=(n-1)\pi(-1)^n.
    \]
    二者不相等,故 $f$ 在每个整数点都不可微,在其余各点可微.
\end{solution}



\subsection{高阶导数}


\methodsection{(A) 利用已知结论}

\begin{theorem}[\markstar]
    常用的 $n$ 阶导数公式为
    \[
        (x^k)^{(n)}=k(k-1)\cdots(k-n+1)x^{k-n},
        \qquad n\leq k,
    \]
    \[
        (e^x)^{(n)}=e^x,
        \qquad
        (a^x)^{(n)}=a^x(\ln a)^n,
    \]
    \[
        (\ln x)^{(n)}=(-1)^{n-1}\dfrac{(n-1)!}{x^n},
    \]
    \[
        (\sin x)^{(n)}=\sin\left(x+\dfrac{n\pi}{2}\right),
        \qquad
        (\cos x)^{(n)}=\cos\left(x+\dfrac{n\pi}{2}\right).
    \]
\end{theorem}

\begin{remark}
    因此有了以上常见 $n$ 阶导数公式,对于一般不易直接求 $n$ 阶导数的函数,我们优先考虑转化成如上形式函数.
\end{remark}

%例题4.1.8
\begin{example}
    计算 $y^{(n)}$,设
    \begin{enumerate}
        \item $y=\sin^6x+\cos^6x$;
        \item $y=\sin ax\cdot\sin bx$;
        \item $y=\dfrac{x^2+x+1}{x^2-5x+6}$.
    \end{enumerate}
\end{example}
\exampleblank

\begin{solution}
    \begin{enumerate}
        \item 由
              \[
                  \sin^6x+\cos^6x
                  =1-3\sin^2x\cos^2x
                  =\dfrac{5}{8}+\dfrac{3}{8}\cos4x,
              \]
              当 $n\geq1$ 时
              \[
                  y^{(n)}=\dfrac{3}{8}4^n
                  \cos\left(4x+\dfrac{n\pi}{2}\right).
              \]
        \item 积化和差得
              \[
                  y=\dfrac{1}{2}\left[\cos(a-b)x-\cos(a+b)x\right],
              \]
              因而
              \[
                  y^{(n)}=\dfrac{1}{2}\left[
                      (a-b)^n\cos\left((a-b)x+\dfrac{n\pi}{2}\right)
                      -(a+b)^n\cos\left((a+b)x+\dfrac{n\pi}{2}\right)
                      \right].
              \]
        \item 部分分式分解为
              \[
                  y=1-\dfrac{7}{x-2}+\dfrac{13}{x-3}.
              \]
              故当 $n\geq1$ 时
              \[
                  y^{(n)}=(-1)^n n!\left[
                      -\dfrac{7}{(x-2)^{n+1}}+\dfrac{13}{(x-3)^{n+1}}
                      \right].
              \]
    \end{enumerate}
\end{solution}



\methodsection{(B) 利用 Leibniz 公式}

%例题4.1.9
\begin{example}
    设 $f_1(x),f_2(x)$ 在 $x_0$ 及其附近有定义,在 $x_0$ 处有直到 $n$ 阶导数,记
    \[
        N(f)=\sum\limits_{k=0}^n\dfrac{1}{k!}|f^{(k)}(x_0)|.
    \]
    试证明:
    \[
        N(f_1f_2)\leq N(f_1)N(f_2).
    \]
\end{example}
\exampleblank

\begin{solution}
    由 Leibniz 公式,对 $0\leq j\leq n$ 有
    \[
        |(f_1f_2)^{(j)}(x_0)|
        \leq\sum\limits_{k=0}^{j}C_{j}^{k}
        |f_1^{(k)}(x_0)|\,|f_2^{(j-k)}(x_0)|.
    \]
    因此
    \[
        \begin{aligned}
            N(f_1f_2)
             & \leq\sum\limits_{j=0}^{n}\sum\limits_{k=0}^{j}
            \dfrac{|f_1^{(k)}(x_0)|}{k!}
            \dfrac{|f_2^{(j-k)}(x_0)|}{(j-k)!}                \\
             & =\sum\limits_{\substack{k,l\geq0               \\k+l\leq n}}
            \dfrac{|f_1^{(k)}(x_0)|}{k!}
            \dfrac{|f_2^{(l)}(x_0)|}{l!}                      \\
             & \leq N(f_1)N(f_2).
        \end{aligned}
    \]
\end{solution}

%例题4.1.10
\begin{example}
    试证:
    \[
        (a^2+b^2)^{\frac{n}{2}}e^{ax}\sin(bx+n\varphi)
        =\sum\limits_{i=0}^n C_{n}^{i}a^{n-i}b^i e^{ax}
        \sin\left(bx+\dfrac{\pi i}{2}\right),
    \]
    其中
    \[
        \varphi=\arctan\dfrac{b}{a}.
    \]
\end{example}
\exampleblank

\begin{solution}
    令 $r=\sqrt{a^2+b^2}$,并选取角 $\varphi$ 使
    \[
        \cos\varphi=\dfrac{a}{r},
        \qquad
        \sin\varphi=\dfrac{b}{r}.
    \]
    直接归纳可得
    \[
        \dfrac{\mathrm{d}^n}{\mathrm{d}x^n}
        \left(e^{ax}\sin bx\right)
        =r^n e^{ax}\sin(bx+n\varphi).
    \]
    另一方面,对 $e^{ax}\sin bx$ 使用 Leibniz 公式,得
    \[
        \begin{aligned}
            \dfrac{\mathrm{d}^n}{\mathrm{d}x^n}
            \left(e^{ax}\sin bx\right)
             & =\sum\limits_{i=0}^{n}C_{n}^{i}
            a^{n-i}e^{ax}b^i
            \sin\left(bx+\dfrac{\pi i}{2}\right).
        \end{aligned}
    \]
    比较两式,并代入 $r^n=(a^2+b^2)^{\frac{n}{2}}$,即得所证恒等式.
\end{solution}



\methodsection{(C) 利用数学归纳法}

当高阶导数不易求出时,可考虑先求出前几阶导数,观察得出规律,进而利用数学归纳法证明.

%例题4.1.11
\begin{example}[\markstar]
    证明:若函数
    \[
        f(x)=
        \begin{cases}
            e^{-\frac{1}{x^2}}, & x\neq0, \\
            0,                  & x=0,
        \end{cases}
    \]
    则在 $x=0$ 处,
    \[
        f^{(n)}(0)=0,
        \qquad n=1,2,\ldots.
    \]
\end{example}
\exampleblank

\begin{solution}
    当 $x\neq0$ 时,对任意 $n\geq0$,反复求导可将 $f^{(n)}$ 写成
    \[
        f^{(n)}(x)=P_n\left(\dfrac{1}{x}\right)e^{-\frac{1}{x^2}},
    \]
    其中 $P_n$ 是多项式.这由归纳立即得到:若对 $n$ 成立,则求导后仍是 $\dfrac{1}{x}$ 的多项式乘以 $e^{-\frac{1}{x^2}}$.

    对任意正整数 $m$,都有
    \[
        \lim\limits_{x\to0}|x|^{-m}e^{-\frac{1}{x^2}}=0.
    \]
    因而若已知 $f^{(n)}(0)=0$,则
    \[
        f^{(n+1)}(0)
        =\lim\limits_{h\to0}\dfrac{f^{(n)}(h)-f^{(n)}(0)}{h}
        =\lim\limits_{h\to0}
        \dfrac{P_n\left(\frac{1}{h}\right)e^{-\frac{1}{h^2}}}{h}=0.
    \]
    由 $f(0)=0$ 归纳得 $f^{(n)}(0)=0$,$n=1,2,\ldots$.
\end{solution}

%例题4.1.12
\begin{example}[推广的 Leibniz]
    设 $u_1(x),u_2(x),\ldots,u_k(x)$ 都是 $x$ 的 $n$ 次可微函数,试证:
    \[
        (u_1u_2\cdots u_k)^{(n)}
        =\sum\limits_{r_1+\cdots+r_k=n}
        \dfrac{n!}{r_1!r_2!\cdots r_k!}
        u_1^{(r_1)}\cdots u_k^{(r_k)}.
    \]
\end{example}
\exampleblank

\begin{solution}
    对 $k$ 作归纳.$k=2$ 时即为 Leibniz 公式.假设结论对 $k-1$ 个因子成立,令
    $v=u_1u_2\cdots u_{k-1}$,则
    \[
        \begin{aligned}
            (vu_k)^{(n)}
             & =\sum\limits_{s=0}^{n}C_{n}^{s}v^{(s)}u_k^{(n-s)} \\
             & =\sum\limits_{s=0}^{n}C_{n}^{s}
            \sum\limits_{r_1+\cdots+r_{k-1}=s}
            \dfrac{s!}{r_1!\cdots r_{k-1}!}
            u_1^{(r_1)}\cdots u_{k-1}^{(r_{k-1})}u_k^{(n-s)}.
        \end{aligned}
    \]
    令 $r_k=n-s$,并利用
    \[
        C_n^s\dfrac{s!}{r_1!\cdots r_{k-1}!}
        =\dfrac{n!}{r_1!\cdots r_k!},
    \]
    即得题设公式.
\end{solution}

%例题4.1.13
\begin{example}[\markstar]
    $f(x)$ 有任意阶导数,$f'(x)=(f(x))^2$,求 $f^{(n)}(x)\ (n>2)$.
\end{example}
\exampleblank

\begin{solution}
    对 $n$ 作归纳.当 $n=1$ 时,题设已给出
    \[
        f'(x)=f(x)^2=1!f(x)^2.
    \]
    若 $f^{(n)}(x)=n!f(x)^{n+1}$,则
    \[
        \begin{aligned}
            f^{(n+1)}(x)
             & =n!(n+1)f(x)^n f'(x) \\
             & =(n+1)!f(x)^{n+2}.
        \end{aligned}
    \]
    因此对所有 $n\geq1$,
    \[
        f^{(n)}(x)=n!f(x)^{n+1}.
    \]
\end{solution}



\methodsection{(D) 利用递推公式}

当更高阶导数可以用低阶导数表示且易于推广到一般情况时,可考虑利用递推关系求更高阶导数.

%例题4.1.14
\begin{example}
    设 $f(x)=(\arcsin x)^2$,求 $f^{(n)}(0)$.
\end{example}
\exampleblank

\begin{solution}
    令 $y=(\arcsin x)^2$.直接求导可得微分方程
    \[
        (1-x^2)y''-xy'=2.
    \]
    记 $a_n=y^{(n)}(0)$.对上式比较 $x^m$ 的系数,或对方程求 $m$ 次导后令 $x=0$,得
    \[
        a_{m+2}=m^2a_m,
        \qquad m\geq1.
    \]
    又 $a_0=a_1=0$,$a_2=2$.因而所有奇数阶导数都为 $0$,而对 $k\geq1$,
    \[
        a_{2k}=2\cdot2^2\cdot4^2\cdots(2k-2)^2
        =2^{2k-1}[(k-1)!]^2.
    \]
    即
    \[
        f^{(n)}(0)=
        \begin{cases}
            0, & n\text{ 为奇数},  \\
            2^{n-1}\left[\left(\dfrac{n}{2}-1\right)!\right]^2,
               & n\text{ 为正偶数}.
        \end{cases}
    \]
\end{solution}

%例题4.1.15
\begin{example}
    证明 Legendre 多项式
    \[
        p_n(x)=\dfrac{1}{2^n n!}\left[(x^2-1)^n\right]^{(n)},
        \qquad n=0,1,2,\ldots,
    \]
    满足方程
    \[
        (1-x^2)p_n''(x)-2xp_n'(x)+n(n+1)p_n(x)=0.
    \]
\end{example}
\exampleblank

\begin{solution}
    令 $u(x)=(x^2-1)^n$,则
    \[
        (x^2-1)u'(x)=2nx\,u(x).
    \]
    对上式两端同时求 $n+1$ 次导.由于 $x^2-1$ 与 $x$ 的高阶导数很快为零,Leibniz 公式给出
    \[
        (x^2-1)u^{(n+2)}+2x u^{(n+1)}-n(n+1)u^{(n)}=0.
    \]
    由定义
    \[
        p_n=\dfrac{u^{(n)}}{2^n n!},
        \qquad
        p_n'=\dfrac{u^{(n+1)}}{2^n n!},
        \qquad
        p_n''=\dfrac{u^{(n+2)}}{2^n n!}.
    \]
    将上式除以 $2^n n!$ 并整体乘以 $-1$,即得
    \[
        (1-x^2)p_n''-2xp_n'+n(n+1)p_n=0.
    \]
\end{solution}



\methodsection{(E) 利用 Taylor 展开}

%例题4.1.16
\begin{example}
    求 $f(x)=\arctan x$ 在 $x=0$ 处的各阶导数.
\end{example}
\exampleblank

\begin{solution}
    在 $|x|<1$ 内,
    \[
        \arctan x
        =\sum\limits_{m=0}^{\infty}(-1)^m\dfrac{x^{2m+1}}{2m+1}.
    \]
    因而偶数阶导数在原点都为零,而
    \[
        f^{(2m+1)}(0)
        =(2m+1)!\dfrac{(-1)^m}{2m+1}
        =(-1)^m(2m)!,
        \qquad m=0,1,2,\ldots.
    \]
\end{solution}

%例题4.1.17
\begin{example}
    设
    \[
        f(x)=
        \begin{cases}
            \dfrac{\sin x}{x}, & x\neq0, \\
            1,                 & x=0,
        \end{cases}
    \]
    求 $f^{(k)}(0)$.
\end{example}
\exampleblank

\begin{solution}
    在原点附近,
    \[
        f(x)=\dfrac{\sin x}{x}
        =\sum\limits_{m=0}^{\infty}(-1)^m\dfrac{x^{2m}}{(2m+1)!}.
    \]
    故当 $k$ 为奇数时 $f^{(k)}(0)=0$;当 $k=2m$ 时,
    \[
        f^{(2m)}(0)
        =(2m)!\dfrac{(-1)^m}{(2m+1)!}
        =\dfrac{(-1)^m}{2m+1}.
    \]
    即
    \[
        f^{(k)}(0)=
        \begin{cases}
            0,                               & k\text{ 为奇数}, \\
            \dfrac{(-1)^{\frac{k}{2}}}{k+1}, & k\text{ 为偶数}.
        \end{cases}
    \]
\end{solution}



\subsection{利用导数求和}

%例题4.1.18
\begin{example}
    试求下列极限:
    \begin{enumerate}
        \item $\displaystyle
                  I=\lim\limits_{x\to0}\dfrac{1}{x}\sum\limits_{i=1}^k f\left(\dfrac{x}{i}\right)$,
              已知 $f'(0)$ 存在,且 $f(0)=0,k\in\mathbb{N}$;
        \item $\displaystyle
                  I=\lim\limits_{n\to\infty}n\left[\sum\limits_{i=1}^k f\left(a+\dfrac{i}{n}\right)-kf(a)\right]$,
              已知 $f'(a)$ 存在;
        \item $\displaystyle
                  I=\lim\limits_{n\to\infty}
                  \dfrac{\left(a+\dfrac{1}{n}\right)^n\left(a+\dfrac{2}{n}\right)^n\cdots\left(a+\dfrac{k}{n}\right)^n}{a^{nk}}$,
              $k\in\mathbb{N}$.
    \end{enumerate}
\end{example}
\exampleblank

\begin{solution}
    \begin{enumerate}
        \item 由 $f(x)=f'(0)x+o(x)$,
              \[
                  \dfrac{1}{x}\sum\limits_{i=1}^{k}f\left(\dfrac{x}{i}\right)
                  =f'(0)\sum\limits_{i=1}^{k}\dfrac{1}{i}+o(1),
              \]
              故 $I=f'(0)\displaystyle\sum\limits_{i=1}^{k}\dfrac{1}{i}$.
        \item 对固定的 $i$,
              \[
                  f\left(a+\dfrac{i}{n}\right)-f(a)
                  =\dfrac{i}{n}f'(a)+o\left(\dfrac{1}{n}\right).
              \]
              对 $i=1,\ldots,k$ 求和得
              \[
                  I=f'(a)\sum\limits_{i=1}^{k}i
                  =\dfrac{k(k+1)}{2}f'(a).
              \]
        \item 假设 $a>0$.取对数后,
              \[
                  \begin{aligned}
                      \ln I_n
                       & =n\sum\limits_{i=1}^{k}
                      \ln\left(1+\dfrac{i}{an}\right)                      \\
                       & \longrightarrow\dfrac{1}{a}\sum\limits_{i=1}^{k}i
                      =\dfrac{k(k+1)}{2a}.
                  \end{aligned}
              \]
              因而
              \[
                  I=\exp\left(\dfrac{k(k+1)}{2a}\right).
              \]
    \end{enumerate}
\end{solution}

%例题4.1.19
\begin{example}
    试求下列和式:
    \begin{enumerate}
        \item $\displaystyle I_n=\sum\limits_{k=0}^n ke^{kx}$;
        \item $\displaystyle \theta_n(x)=\sum\limits_{k=1}^n k^2x^{k-1}$;
        \item $\displaystyle I_n=\sum\limits_{k=1}^n C_{n}^{k}k^2$;
        \item $\displaystyle I_n=\sum\limits_{k=0}^{2n}(-1)^k C_{2n}^{k}k^n$;
        \item $\displaystyle S_n(x)=\sum\limits_{k=1}^n\dfrac{1}{2^k}\tan\left(\dfrac{x}{2^k}\right)$.
    \end{enumerate}
\end{example}
\exampleblank

\begin{solution}
    \begin{enumerate}
        \item 令 $q=e^x$.对等比和求导,得
              \[
                  I_n=\sum\limits_{k=0}^{n}kq^k
                  =\dfrac{q-(n+1)q^{n+1}+nq^{n+2}}{(1-q)^2}.
              \]
        \item 由
              \[
                  \sum\limits_{k=1}^{n}kx^k
                  =\dfrac{x-(n+1)x^{n+1}+nx^{n+2}}{(1-x)^2},
              \]
              两端对 $x$ 求导得,当 $x\neq1$ 时
              \[
                  \theta_n(x)=
                  \dfrac{1+x-(n+1)^2x^n+(2n^2+2n-1)x^{n+1}-n^2x^{n+2}}
                  {(1-x)^3}.
              \]
              当 $x=1$ 时,
              \[
                  \theta_n(1)=\sum\limits_{k=1}^{n}k^2
                  =\dfrac{n(n+1)(2n+1)}{6}.
              \]
        \item 对 $(1+x)^n$ 求导一次,两次,并在 $x=1$ 处取值,得
              \[
                  \begin{aligned}
                      \sum\limits_{k=1}^{n}C_n^k k^2
                       & =\sum\limits_{k=1}^{n}C_n^k k(k-1)
                      +\sum\limits_{k=1}^{n}C_n^k k         \\
                       & =n(n-1)2^{n-2}+n2^{n-1}            \\
                       & =n(n+1)2^{n-2}.
                  \end{aligned}
              \]
        \item 该和是多项式 $P(k)=k^n$ 的 $2n$ 阶有限差(相差一个符号).由于 $\deg P=n<2n$,其 $2n$ 阶有限差为零,故
              \[
                  I_n=0.
              \]
        \item 使用恒等式
              \[
                  \tan u=\cot u-2\cot2u.
              \]
              于是
              \[
                  \begin{aligned}
                      S_n(x)
                       & =\sum\limits_{k=1}^{n}\left[
                          \dfrac{1}{2^k}\cot\dfrac{x}{2^k}
                          -\dfrac{1}{2^{k-1}}\cot\dfrac{x}{2^{k-1}}
                      \right]                                      \\
                       & =\dfrac{1}{2^n}\cot\dfrac{x}{2^n}-\cot x.
                  \end{aligned}
              \]
    \end{enumerate}
\end{solution}

%例题4.1.20
\begin{example}
    设 $f_1(x),f_2(x),\ldots,f_n(x)$ 在 $x=x_0$ 处可导,且 $f_k(x_0)\neq0\ (k=1,2,\ldots,n)$,
    \[
        F(x)=\prod\limits_{k=1}^n f_k(x),
    \]
    则
    \[
        \dfrac{F'(x_0)}{F(x_0)}
        =\sum\limits_{k=1}^n\dfrac{f_k'(x_0)}{f_k(x_0)}.
    \]
\end{example}
\exampleblank

\begin{solution}
    由有限个函数乘积的求导法则,
    \[
        F'(x_0)=\sum\limits_{j=1}^{n}f_j'(x_0)
        \prod\limits_{\substack{1\leq k\leq n\\k\neq j}}f_k(x_0).
    \]
    因为所有 $f_k(x_0)$ 都非零,两端除以
    $F(x_0)=\displaystyle\prod\limits_{k=1}^{n}f_k(x_0)$,即得
    \[
        \dfrac{F'(x_0)}{F(x_0)}
        =\sum\limits_{j=1}^{n}\dfrac{f_j'(x_0)}{f_j(x_0)}.
    \]
\end{solution}



\subsection{利用导数证恒等式}

%例题4.1.21
\begin{example}
    证明:
    \[
        2\arctan x+\arcsin\dfrac{2x}{1+x^2}=\pi,
        \qquad 1<x<\infty.
    \]
\end{example}
\exampleblank

\begin{solution}
    令
    \[
        F(x)=2\arctan x+\arcsin\dfrac{2x}{1+x^2}.
    \]
    当 $x>1$ 时,
    \[
        \sqrt{1-\left(\dfrac{2x}{1+x^2}\right)^2}
        =\dfrac{x^2-1}{1+x^2}.
    \]
    因而
    \[
        F'(x)=\dfrac{2}{1+x^2}
        +\dfrac{2(1-x^2)}{(1+x^2)^2}
        \dfrac{1+x^2}{x^2-1}=0.
    \]
    故 $F$ 在 $(1,+\infty)$ 上为常数.令 $x\to1^+$,得
    \[
        F(x)\longrightarrow2\cdot\dfrac{\pi}{4}+\dfrac{\pi}{2}=\pi.
    \]
    因此题设恒等式成立.
\end{solution}

%例题4.1.22
\begin{example}
    设 $f(x),g(x)$ 在 $(-\infty,+\infty)$ 上均不是常数,可导且有
    \[
        \begin{cases}
            f(x+y)=f(x)f(y)-g(x)g(y), \\
            g(x+y)=f(x)g(y)+g(x)f(y),
        \end{cases}
    \]
    其中 $x,y\in(-\infty,+\infty)$,若 $f'(0)=0$,则
    \[
        f^2(x)+g^2(x)=1.
    \]
\end{example}
\exampleblank

\begin{solution}
    在加法公式中令 $y=0$.由于 $f,g$ 不同时为常数,可得
    \[
        f(0)=1,
        \qquad
        g(0)=0.
    \]
    对两个加法公式关于 $y$ 求导,再令 $y=0$.记 $c=g'(0)$,并使用 $f'(0)=0$,得
    \[
        f'(x)=-c g(x),
        \qquad
        g'(x)=c f(x).
    \]
    因此
    \[
        \dfrac{\mathrm{d}}{\mathrm{d}x}\left[f(x)^2+g(x)^2\right]
        =2ff'+2gg'=0.
    \]
    该和为常数,并且在 $x=0$ 处等于 $1$,故
    \[
        f(x)^2+g(x)^2=1.
    \]
\end{solution}

%例题4.1.23
\begin{example}
    设 $f(x),g(x)$ 在 $(a,b)$ 上均不是常数,若有
    \[
        f(x)+g(x)\neq0,
        \qquad
        f(x)g'(x)-f'(x)g(x)=0,
        \qquad x\in(a,b),
    \]
    则 $g(x)\neq0\ (x\in(a,b))$ 且
    \[
        \dfrac{f(x)}{g(x)}\equiv C\text{(常数)}.
    \]
\end{example}
\exampleblank

\begin{solution}
    由 $f+g\neq0$,函数
    \[
        H(x)=\dfrac{f(x)}{f(x)+g(x)}
    \]
    在 $(a,b)$ 上有定义.题设等式给出
    \[
        H'(x)=\dfrac{f'(x)g(x)-f(x)g'(x)}{[f(x)+g(x)]^2}=0.
    \]
    故 $H(x)\equiv A$ 为常数,从而
    \[
        \dfrac{g(x)}{f(x)+g(x)}\equiv1-A.
    \]
    若某点 $g(x_0)=0$,则 $1-A=0$,上式迫使 $g\equiv0$,与 $g$ 不是常数矛盾.因而 $g(x)\neq0$.又 $1-A\neq0$,所以
    \[
        \dfrac{f(x)}{g(x)}
        =\dfrac{A}{1-A}
    \]
    为常数.
\end{solution}



\subsection{导函数的性质}

导函数有两个重要基本定理---Darboux 定理和单侧导数极限定理,进而可得导数的两个重要特性:(1) 介值性;(2) 无第一类间断点.

\begin{theorem}[\markstar\ Darboux 定理]
    设 $f$ 在区间 $I$ 上可微,则 $f'$ 具有介值性.
\end{theorem}

\begin{corollary}[\markstar]
    单调的导函数 $f'$ 必连续.
\end{corollary}

\begin{theorem}[\markstar\ 单侧导数极限定理]
    设在区间 $\lbrack a,b)$ 上有定义的函数 $f$ 在 $(a,b)$ 上可微,又在点 $a$ 右连续,若导函数 $f'(x)$ 在点 $a$ 存在右侧极限
    \[
        f'(a+0)=A,
    \]
    则 $f$ 在点 $a$ 也一定存在右侧导数 $f'_+(a)$,且成立
    \[
        f'_+(a)=f'(a+0)=A.
    \]
    即 $f'(x)$ 在点 $a$ 右连续.
\end{theorem}

\begin{corollary}[\markstar]
    设 $f$ 在 $(a,b)$ 上可微,则导函数 $f'(x)$ 不存在第一类间断点.
\end{corollary}

%例题4.1.24
\begin{example}
    设函数 $f(x)$ 在区间 $(-\infty,+\infty)$ 上二次可微,且有界,试证存在点 $x_0\in(-\infty,+\infty)$,使得 $f'(x_0)=0$.
\end{example}
\exampleblank

\begin{solution}
    本题按当前题面不成立.取
    \[
        f(x)=\arctan x.
    \]
    则 $f$ 在实轴上二次可微且有界,但
    \[
        f'(x)=\dfrac{1}{1+x^2}>0
    \]
    对所有 $x\in\mathbb{R}$ 成立,因而不存在 $f'(x_0)=0$ 的点.原题需要补充条件或核对结论.
\end{solution}

%例题4.1.25
\begin{example}
    设函数 $f$ 在 $(0,1)$ 有界,可导,但 $\displaystyle\lim\limits_{x\to0^+}f(x)$ 不存在,试证存在数列 $x_n\to0^+\ (n\to\infty)$,使得
    \[
        \lim\limits_{n\to+\infty}f'(x_n)=0.
    \]
\end{example}
\exampleblank

\begin{solution}
    反证.若不存在所求数列,则存在 $\varepsilon>0$ 与 $\delta>0$,使
    \[
        |f'(x)|\geq\varepsilon,
        \qquad 0<x<\delta.
    \]
    导函数具有 Darboux 介值性,因而 $f'$ 在 $(0,\delta)$ 上不能变号.若 $f'\geq\varepsilon$,则对 $0<x<y<\delta$,中值定理给出
    \[
        f(y)-f(x)\geq\varepsilon(y-x),
    \]
    故 $f$ 随 $x\to0^+$ 单调且有界,从而存在有限右极限;$f'\leq-\varepsilon$ 时同理.这与题设矛盾.

    因此对每个 $n$,都能取 $0<x_n<\dfrac{1}{n}$ 使 $|f'(x_n)|<\dfrac{1}{n}$.于是 $x_n\to0^+$ 且 $f'(x_n)\to0$.
\end{solution}

%例题4.1.26
\begin{example}
    \begin{enumerate}
        \item 设 $f(x)$ 在 $(a,b)$ 上可导,且有 $x_i\in(a,b),\lambda_i>0\ (i=1,2,\ldots,n)$,
              \[
                  \sum\limits_{i=1}^n\lambda_i=1,
              \]
              则存在 $\xi\in(a,b)$,使得
              \[
                  \sum\limits_{i=1}^n\lambda_i f'(x_i)=f'(\xi);
              \]
        \item 设 $f(x)$ 在 $(a,b)$ 上可导,且有 $x_i<y_i$,$x_i,y_i\in(a,b)\ (i=1,2,\ldots,n)$,则存在 $\xi\in(a,b)$,使得
              \[
                  \sum\limits_{i=1}^n[f(y_i)-f(x_i)]
                  =f'(\xi)\sum\limits_{i=1}^n(y_i-x_i).
              \]
    \end{enumerate}
\end{example}
\exampleblank

\begin{solution}
    \begin{enumerate}
        \item 数 $\displaystyle\sum\limits_{i=1}^{n}\lambda_i f'(x_i)$ 是 $f'(x_i)$ 的凸组合,因而位于这些值的最小值与最大值之间.由导函数的 Darboux 介值性,存在 $\xi\in(a,b)$ 使
              \[
                  f'(\xi)=\sum\limits_{i=1}^{n}\lambda_i f'(x_i).
              \]
        \item 由 Lagrange 中值定理,对每个 $i$ 存在 $c_i\in(x_i,y_i)$,使
              \[
                  f(y_i)-f(x_i)=f'(c_i)(y_i-x_i).
              \]
              令
              \[
                  \lambda_i=\dfrac{y_i-x_i}{\displaystyle\sum\limits_{j=1}^{n}(y_j-x_j)},
              \]
              则 $\lambda_i>0$ 且其和为 $1$.对 $f'(c_i)$ 应用第 (1) 问,即得题设等式.
    \end{enumerate}
\end{solution}

%例题4.1.27
\begin{example}
    设 $f(x)$ 在 $\lbrack a,b\rbrack$ 上二次可导,若 $f''(x)\neq0\ (a\leq x\leq b)$,则对任意 $\xi\in(a,b)$,存在 $\xi'$ 满足 $a<\xi<\xi'\leq b$,使得
    \[
        f'(\xi)=\dfrac{f(\xi')-f(a)}{\xi'-a}.
    \]
\end{example}
\exampleblank

\begin{solution}
    本题按当前题面不成立.取 $\lbrack a,b\rbrack=\lbrack0,1\rbrack$ 与
    \[
        f(x)=e^x.
    \]
    则 $f''(x)=e^x\neq0$.若取 $\xi$ 充分接近 $1$,例如 $\xi=\dfrac{9}{10}$,则
    \[
        f'(\xi)=e^{\frac{9}{10}}>e-1.
    \]
    但对任意 $\xi'\in(\xi,1\rbrack$,
    \[
        \dfrac{f(\xi')-f(0)}{\xi'}
        =\dfrac{e^{\xi'}-1}{\xi'}
        \leq e-1<f'(\xi).
    \]
    因而不存在题设的 $\xi'$.原题中``对任意 $\xi$''或割线端点条件需要核对.
\end{solution}



\subsection{广义导数 *}

\begin{definition}
    $\Delta_hf(x)=f(x+h)-f(x-h)$ 称为 $f(x)$ 在点 $x$ 处的对称差,如下极限(若存在)称为 $f(x)$ 在点 $x$ 处的广义导数或对称导数,记作
    \[
        f^{[\prime]}(x)
        =\lim\limits_{h\to0^+}\dfrac{\Delta_hf(x)}{2h}
        =\lim\limits_{h\to0^+}\dfrac{f(x+h)-f(x-h)}{2h}.
    \]
\end{definition}

\begin{theorem}[\markstar]
    设 $f$ 在点 $x$ 处可导,则 $f^{[\prime]}$ 存在且
    \[
        f^{[\prime]}(x)=f'(x).
    \]
\end{theorem}

\begin{remark}
    该定理逆命题不成立,例如
    \[
        f(x)=
        \begin{cases}
            x\sin\dfrac{1}{x}, & x\neq0, \\
            0,                 & x=0,
        \end{cases}
    \]
    或 $f(x)=|x|$.
\end{remark}

\begin{definition}
    如下的极限存在,称为 $f(x)$ 在点 $x$ 处有二阶广义导数(或称二阶对称导数),记作
    \[
        \begin{aligned}
            f^{[\prime\prime]}(x)
             & =f^{[2]}(x)                                            \\
             & \overset{\triangle}{=}
            \lim\limits_{h\to0^+}\dfrac{\Delta_h\Delta_hf(x)}{(2h)^2} \\
             & =\lim\limits_{h\to0^+}
            \dfrac{\Delta_h[f(x+h)-f(x-h)]}{(2h)^2}                   \\
             & =\lim\limits_{h\to0^+}
            \dfrac{f(x+2h)-2f(x)-f(x-2h)}{4h^2}.
        \end{aligned}
    \]
\end{definition}

\begin{theorem}[\markstar]
    在某点 $x$,若 $f(x)$ 有二阶导数,则二阶广义导数也存在,且值相等,即
    \[
        f''(x)=f^{[2]}(x).
    \]
\end{theorem}

\begin{remark}
    逆命题同样不成立,例如
    \[
        f(x)=
        \begin{cases}
            x^2\sin\dfrac{1}{x}, & x\neq0, \\
            0,                   & x=0,
        \end{cases}
    \]
    或 $f(x)=x|x|$.
\end{remark}

%例题4.1.28
\begin{example}[\markstar\ Schwarz 定理]
    若 $f(x)$ 在 $\lbrack a,b\rbrack$ 上连续,$f(x)$ 的广义二阶导数 $f^{[2]}(x)$ 存在且恒为零,则 $f(x)$ 是一次线性函数,即 $f(x)=ax+b$.
\end{example}
\exampleblank

\begin{solution}
    记 $L(x)$ 为连接 $(a,f(a))$ 与 $(b,f(b))$ 的一次函数,并令
    \[
        u(x)=f(x)-L(x).
    \]
    则 $u(a)=u(b)=0$,且 $u^{[2]}(x)=0$.任取 $\varepsilon>0$,令
    \[
        u_{\varepsilon}(x)=u(x)+\varepsilon x^2.
    \]
    按照二阶对称导数的定义,有
    \[
        u_{\varepsilon}^{[2]}(x)=2\varepsilon>0.
    \]
    若 $u_{\varepsilon}$ 在内点 $x_0$ 取得最大值,则对充分小的 $h>0$,
    \[
        u_{\varepsilon}(x_0+2h)-2u_{\varepsilon}(x_0)
        +u_{\varepsilon}(x_0-2h)\leq0,
    \]
    从而 $u_{\varepsilon}^{[2]}(x_0)\leq0$,矛盾.因此 $u_{\varepsilon}$ 的最大值只能在端点取得,于是
    \[
        u(x)+\varepsilon x^2
        \leq\max\{\varepsilon a^2,\varepsilon b^2\}.
    \]
    令 $\varepsilon\to0^+$,得 $u(x)\leq0$.对 $-u$ 作同样讨论,得 $u(x)\geq0$,故 $u\equiv0$,即 $f=L$ 为一次函数.

    需要说明的是,留白版当前定义式的分子末项写成了 $-f(x-2h)$;上述定理所需的标准二阶对称差分应为 $+f(x-2h)$,该处符号需要对照原始资料核对.
\end{solution}

%例题4.1.29
\begin{example}
    若 $f(x)$ 在 $\lbrack a,b\rbrack$ 上连续,$f(a)<f(b)$,又设对一切 $x\in(a,b)$,
    \[
        \lim\limits_{t\to0}\dfrac{f(x+t)-f(x-t)}{t}
    \]
    存在,用 $g(x)$ 表示这一极限值,则存在 $c\in(a,b)$,使得 $g(c)\geq0$.
\end{example}
\exampleblank

\begin{solution}
    这里使用广义中值定理:若连续函数 $F$ 在 $\lbrack\alpha,\beta\rbrack$ 上满足 $F(\alpha)<F(\beta)$,且其对称导数
    \[
        D_sF(x)=\lim\limits_{t\to0}\dfrac{F(x+t)-F(x-t)}{t}
    \]
    在每个内点存在,则至少存在一点 $c\in(\alpha,\beta)$ 使 $D_sF(c)\geq0$.该结论是普通中值定理在对称导数情形下的 Schwarz 中值定理.

    本题中 $F=f$,题设恰好给出 $f(a)<f(b)$ 以及每个内点处对称导数的存在性.因此由上述广义中值定理,存在 $c\in(a,b)$,使
    \[
        g(c)=D_sf(c)\geq0.
    \]
\end{solution}


\newpage

\section{微分中值定理}


\subsection{构造辅助函数(中值问题)}

\noindent\textbf{(\markstar)} 看到 $f'(\xi)+f(\xi)g(\xi)$,应想到其原函数为
$\displaystyle f(x)e^{\int g(x)\,\mathrm{d}x}$,这是因为
\[
    \left[f(x)e^{\int g(x)\,\mathrm{d}x}\right]'
    =[f'(x)+f(x)g(x)]e^{\int g(x)\,\mathrm{d}x}.
\]
具体地,有如下常见体现:
\[
    mf(\xi)+nf'(\xi)
    \longleftrightarrow
    f(x)e^{\frac{m}{n}x}.
\]
\[
    mf(\xi)+n\xi f'(\xi)
    \longleftrightarrow
    x^m f^n(x).
\]
\[
    mf(\xi)-n\xi f'(\xi)
    \longleftrightarrow
    \dfrac{f^n(x)}{x^m}
    \quad\text{或}\quad
    \dfrac{x^m}{f^n(x)}.
\]
\[
    nf'(\xi)f(1-\xi)-mf(\xi)f'(1-\xi)
    \longleftrightarrow
    f^n(x)f^m(1-x).
\]
\[
    mf'(\xi)g(\xi)+nf(\xi)g'(\xi)
    \longleftrightarrow
    f^m(x)g^n(x).
\]
\[
    mf'(\xi)g(\xi)-nf(\xi)g'(\xi)
    \longleftrightarrow
    \dfrac{f^m(x)}{g^n(x)}.
\]
\[
    f(\xi)g''(\xi)-g(\xi)f''(\xi)
    \longleftrightarrow
    f'(x)g(x)-f(x)g'(x).
\]

%例题4.2.1
\begin{example}
    设 $f\in C(\lbrack a,b\rbrack)$,$g\in C(\lbrack a,b\rbrack)$,且均在 $(a,b)$ 上可导,若 $f(a)=f(b)=0$,则存在 $\xi\in(a,b)$,使得
    \[
        f'(\xi)+g'(\xi)f(\xi)=0.
    \]
\end{example}
\exampleblank

\begin{solution}
    构造辅助函数
    \[
        F(x)=f(x)e^{g(x)}.
    \]
    由 $f(a)=f(b)=0$,有 $F(a)=F(b)=0$.根据 Rolle 定理,存在 $\xi\in(a,b)$ 使 $F'(\xi)=0$.又
    \[
        F'(x)=e^{g(x)}\left[f'(x)+g'(x)f(x)\right].
    \]
    由于 $e^{g(\xi)}>0$,必有
    \[
        f'(\xi)+g'(\xi)f(\xi)=0.
    \]
\end{solution}

%例题4.2.2
\begin{example}
    设 $f(x)$ 在 $\lbrack 0,1\rbrack$ 上二次可导,且 $f(0)=0$,则存在 $\xi\in(0,1)$,使得
    \[
        f'(\xi)-(\xi-1)^2f''(\xi)=0.
    \]
\end{example}
\exampleblank

\begin{solution}
    本题按当前题面不成立.取 $f(x)=x$,则 $f(0)=0$,且 $f$ 在 $\lbrack0,1\rbrack$ 上二次可导,但对每个 $\xi\in(0,1)$ 都有
    \[
        f'(\xi)-(\xi-1)^2f''(\xi)=1\neq0.
    \]
    因此原题还缺少条件,需要核对原始资料.
\end{solution}

%例题4.2.3
\begin{example}
    设 $f(x),g(x)$ 在 $\lbrack a,b\rbrack$ 二次可导,且 $g''(x)\neq0$,若有
    \[
        f(a)=f(b)=g(a)=g(b)=0,
    \]
    则存在 $\xi\in(a,b)$,使得
    \[
        \dfrac{f(\xi)}{g(\xi)}=\dfrac{f''(\xi)}{g''(\xi)}.
    \]
\end{example}
\exampleblank

\begin{solution}
    由 $g''$ 的 Darboux 性质及 $g''(x)\neq0$,$g''$ 在 $(a,b)$ 上保持定号.结合 $g(a)=g(b)=0$ 可知 $g(x)\neq0$,$x\in(a,b)$.

    令
    \[
        H(x)=f'(x)g(x)-f(x)g'(x).
    \]
    由端点条件,$H(a)=H(b)=0$.Rolle 定理给出某个 $\xi\in(a,b)$,使 $H'(\xi)=0$.而
    \[
        H'(x)=f''(x)g(x)-f(x)g''(x),
    \]
    故
    \[
        f''(\xi)g(\xi)=f(\xi)g''(\xi).
    \]
    由于 $g(\xi)g''(\xi)\neq0$,两端相除即得结论.
\end{solution}

%例题4.2.4
\begin{example}
    设 $f(x),g(x)$ 在 $\lbrack a,b\rbrack$ 上可导,且在 $(a,b)$ 上 $g'(x)\neq0$,则存在 $\xi\in(a,b)$,使得
    \[
        \dfrac{f(a)-f(\xi)}{g(\xi)-g(b)}=\dfrac{f'(\xi)}{g'(\xi)}.
    \]
\end{example}
\exampleblank

\begin{solution}
    令
    \[
        H(x)=[f(a)-f(x)][g(x)-g(b)].
    \]
    则 $H(a)=H(b)=0$.由 Rolle 定理,存在 $\xi\in(a,b)$ 使 $H'(\xi)=0$,即
    \[
        -f'(\xi)[g(\xi)-g(b)]
        +[f(a)-f(\xi)]g'(\xi)=0.
    \]
    又因 $g'$ 不为零,$g$ 严格单调,故 $g(\xi)\neq g(b)$.整理并相除便得题设等式.
\end{solution}

%例题4.2.5
\begin{example}
    设 $f\in C(\lbrack 0,1\rbrack)$,在 $(0,1)$ 上可导且 $f(1)=0$,则存在 $\xi\in(0,1)$ 使得
    \[
        f'(\xi)=\left(1-\dfrac{1}{\xi}\right)f(\xi).
    \]
\end{example}
\exampleblank

\begin{solution}
    定义
    \[
        H(x)=xe^{-x}f(x),
        \qquad 0\leq x\leq1.
    \]
    由 $f$ 的连续性及 $f(1)=0$,有 $H(0)=H(1)=0$.故存在 $\xi\in(0,1)$ 使 $H'(\xi)=0$.计算得
    \[
        H'(x)=xe^{-x}\left[f'(x)-\left(1-\dfrac{1}{x}\right)f(x)\right].
    \]
    因 $\xi e^{-\xi}>0$,括号内的式子在 $x=\xi$ 时为零.
\end{solution}

%例题4.2.6
\begin{example}
    设 $f\in C(\lbrack a,b\rbrack)$,在 $(a,b)$ 上可导,若有
    \[
        f(a)f(b)>0,
        \qquad
        f(a)f\left(\dfrac{a+b}{2}\right)<0,
    \]
    则对任给的 $\lambda\in(-\infty,+\infty)$,存在 $\xi\in(a,b)$,使得
    \[
        f'(\xi)=\lambda f(\xi).
    \]
\end{example}
\exampleblank

\begin{solution}
    记 $m=\dfrac{a+b}{2}$.由 $f(a)f(m)<0$ 以及 $f(a)f(b)>0$,连续性给出两点
    \[
        \alpha\in(a,m),
        \qquad
        \beta\in(m,b),
    \]
    使 $f(\alpha)=f(\beta)=0$.

    对任意给定的 $\lambda\in\mathbb{R}$,令 $H(x)=e^{-\lambda x}f(x)$.则 $H(\alpha)=H(\beta)=0$,故由 Rolle 定理,存在 $\xi\in(\alpha,\beta)$ 使
    \[
        0=H'(\xi)=e^{-\lambda\xi}[f'(\xi)-\lambda f(\xi)].
    \]
    这就是所需结论.
\end{solution}

%例题4.2.7
\begin{example}
    设 $f(x)$ 在 $\lbrack 0,+\infty)$ 上可导,若
    \[
        f^2(a)-f^2(b)=b^2-a^2,
    \]
    则存在 $\xi\in(0,+\infty)$,使得
    \[
        f'(\xi)=\dfrac{1-\xi^2}{(1+\xi^2)^2}.
    \]
\end{example}
\exampleblank

\begin{solution}
    本题按当前题面不成立.任取 $0<a<b$,令
    \[
        f(x)=a+b-x,
        \qquad x\geq0.
    \]
    则
    \[
        f(a)^2-f(b)^2=b^2-a^2,
    \]
    但 $f'(x)=-1$.另一方面,对 $x\geq0$,
    \[
        \dfrac{1-x^2}{(1+x^2)^2}\geq-\dfrac{1}{8}>-1.
    \]
    故不存在题设中的 $\xi$,原题的条件或结论需要核对.
\end{solution}

%例题4.2.8
\begin{example}
    设 $f(x)$ 在 $\lbrack 0,1\rbrack$ 上二次可导,且 $f(0)=f(1)=0$,则存在 $\xi\in(0,1)$,使得
    \[
        f''(\xi)=\dfrac{2f'(\xi)}{1-\xi}.
    \]
\end{example}
\exampleblank

\begin{solution}
    由 $f(0)=f(1)=0$ 及 Rolle 定理,存在 $c\in(0,1)$ 使 $f'(c)=0$.令
    \[
        H(x)=(1-x)^2f'(x).
    \]
    则 $H(c)=0$,且 $\lim\limits_{x\to1^-}H(x)=0$.将 $H$ 连续补定义到 $x=1$ 后,再用 Rolle 定理,存在 $\xi\in(c,1)$ 使
    \[
        0=H'(\xi)
        =(1-\xi)[(1-\xi)f''(\xi)-2f'(\xi)].
    \]
    由于 $\xi<1$,整理即得结论.
\end{solution}

%例题4.2.9
\begin{example}
    设 $f(x)$ 在 $\lbrack 0,1\rbrack$ 上可导,且 $f(0)=0,f(x)>0\ (0<x<1)$,则存在 $\xi\in(0,1)$,使得
    \[
        \dfrac{2f'(\xi)}{f(\xi)}
        =\dfrac{f'(1-\xi)}{f(1-\xi)}.
    \]
\end{example}
\exampleblank

\begin{solution}
    在 $(0,1)$ 上令 $h(x)=\ln f(x)$,并定义
    \[
        H(x)=2h(x)+h(1-x).
    \]
    由 $f(0)=0$,$f(x)>0$ 及连续性,$H(x)\to-\infty$ 当 $x\to0^+$ 或 $x\to1^-$.因此 $H$ 在某个内点 $\xi\in(0,1)$ 取得最大值.于是
    \[
        0=H'(\xi)
        =2\dfrac{f'(\xi)}{f(\xi)}
        -\dfrac{f'(1-\xi)}{f(1-\xi)},
    \]
    即得所证等式.
\end{solution}



\subsection{待定常数法(中值问题)}

该方法的本质是将需要证的中值设为常数,而后将原式子中的某个常数设为变量,再根据 Rolle 或 Lagrange 中值定理证明中值点的存在性.

%例题4.2.10
\begin{example}
    设 $f(x)$ 在 $\lbrack a,b\rbrack$ 上二次可导,则存在 $\xi\in(a,b)$,使得
    \[
        f(b)-2f\left(\dfrac{a+b}{2}\right)+f(a)
        =\dfrac{(b-a)^2}{4}f''(\xi).
    \]
\end{example}
\exampleblank

\begin{solution}
    记 $m=\dfrac{a+b}{2}$.作通过三点 $(a,f(a))$,$(m,f(m))$,$(b,f(b))$ 的二次插值多项式 $P$.函数 $f-P$ 在 $a,m,b$ 三点为零,连续使用两次 Rolle 定理,存在 $\xi\in(a,b)$ 使
    \[
        f''(\xi)=P''(\xi).
    \]
    由于三点等距,直接计算二次差商得
    \[
        P''(x)=\dfrac{4}{(b-a)^2}
        \left[f(b)-2f(m)+f(a)\right].
    \]
    移项即得题设等式.
\end{solution}

%例题4.2.11
\begin{example}
    设 $f(x),g(x)$ 在 $\lbrack a,b\rbrack$ 上二次可导,则存在 $\xi\in(a,b)$,使得
    \[
        \dfrac{f(b)-f(a)-(b-a)f'(a)}{g(b)-g(a)-(b-a)g'(a)}
        =\dfrac{f''(\xi)}{g''(\xi)}.
    \]
\end{example}
\exampleblank

\begin{solution}
    当前题面缺少保证两个比值有意义的非退化条件.例如取 $g(x)=x$,则
    \[
        g(b)-g(a)-(b-a)g'(a)=0,
        \qquad
        g''(x)=0,
    \]
    题设两端均无定义.

    若补充 $g''(x)\neq0$,则先对 $f(x)-f(a)-(x-a)f'(a)$ 与相应的 $g$ 函数应用 Cauchy 中值定理,再对所得的一阶导数差应用一次 Cauchy 中值定理,即可得到
    \[
        \dfrac{f(b)-f(a)-(b-a)f'(a)}
        {g(b)-g(a)-(b-a)g'(a)}
        =\dfrac{f''(\xi)}{g''(\xi)}.
    \]
    因此原题应补上相应的分母非零条件.
\end{solution}

%例题4.2.12
\begin{example}
    若存在 $f''(0)$,则
    \[
        \lim\limits_{h\to0}
        \dfrac{f(2h)-2f(0)+f(-2h)}{4h^2}
        =f''(0).
    \]
\end{example}
\exampleblank

\begin{solution}
    由 $f''(0)$ 存在,$f$ 在 $0$ 处有二阶 Peano 展开:
    \[
        f(x)=f(0)+f'(0)x+\dfrac{1}{2}f''(0)x^2+o(x^2).
    \]
    分别令 $x=2h$ 与 $x=-2h$,相加并消去一次项,得
    \[
        f(2h)-2f(0)+f(-2h)=4h^2f''(0)+o(h^2).
    \]
    两端除以 $4h^2$ 并令 $h\to0$ 即得结论.
\end{solution}

%例题4.2.13
\begin{example}
    设 $h>0$,$f\in C(\lbrack a-h,a+h\rbrack)$,在 $(a-h,a+h)$ 上可导,则存在 $\theta:0<\theta<1$,使得
    \[
        \dfrac{f(a+h)-f(a-h)}{h}
        =f'(a+\theta h)+f'(a-\theta h).
    \]
\end{example}
\exampleblank

\begin{solution}
    在 $\lbrack0,h\rbrack$ 上令
    \[
        \Phi(t)=f(a+t)-f(a-t).
    \]
    由 Lagrange 中值定理,存在 $t_0\in(0,h)$ 使
    \[
        \dfrac{\Phi(h)-\Phi(0)}{h}=\Phi'(t_0).
    \]
    写成 $t_0=\theta h$,其中 $0<\theta<1$.注意到
    \[
        \Phi'(t)=f'(a+t)+f'(a-t),
    \]
    便得到题设等式.
\end{solution}

%例题4.2.14
\begin{example}
    设 $f(x)$ 在 $\lbrack a,b\rbrack$ 上连续,且在 $(a,b)$ 上二次可导,则对 $x\in(a,b)$,存在 $\xi\in(a,b)$,使得
    \[
        \dfrac{f(x)-f(a)}{x-a}
        -\dfrac{f(b)-f(a)}{b-a}
        =\dfrac{1}{2}(x-b)f''(\xi).
    \]
\end{example}
\exampleblank

\begin{solution}
    令 $L(t)$ 为连接 $(a,f(a))$ 与 $(b,f(b))$ 的一次函数,并取常数 $K$ 使
    \[
        \Phi(t)=f(t)-L(t)-K(t-a)(t-b)
    \]
    还满足 $\Phi(x)=0$.于是 $\Phi(a)=\Phi(x)=\Phi(b)=0$.连续使用两次 Rolle 定理,存在 $\xi\in(a,b)$ 使 $\Phi''(\xi)=0$,故
    \[
        K=\dfrac{1}{2}f''(\xi).
    \]
    另一方面,由 $\Phi(x)=0$,
    \[
        K=\dfrac{f(x)-L(x)}{(x-a)(x-b)}.
    \]
    又
    \[
        \dfrac{f(x)-L(x)}{x-a}
        =\dfrac{f(x)-f(a)}{x-a}
        -\dfrac{f(b)-f(a)}{b-a}.
    \]
    合并以上各式即得结论.
\end{solution}

%例题4.2.15
\begin{example}
    设 $f(x)$ 在 $\lbrack a,c\rbrack$ 上二次可导,$a<b<c$,则存在 $\xi\in(a,c)$,使得
    \[
        \dfrac{f(a)}{(a-b)(a-c)}
        +\dfrac{f(b)}{(b-c)(b-a)}
        +\dfrac{f(c)}{(c-a)(c-b)}
        =\dfrac{1}{2}f''(\xi).
    \]
\end{example}
\exampleblank

\begin{solution}
    设 $P$ 是通过 $(a,f(a))$,$(b,f(b))$,$(c,f(c))$ 的二次插值多项式.由插值余项的 Rolle 定理证明,存在 $\xi\in(a,c)$ 使
    \[
        f''(\xi)=P''(\xi).
    \]
    将 $P$ 写成 Lagrange 形式并求二阶导数,得到
    \[
        \dfrac{1}{2}P''(x)
        =\dfrac{f(a)}{(a-b)(a-c)}
        +\dfrac{f(b)}{(b-c)(b-a)}
        +\dfrac{f(c)}{(c-a)(c-b)}.
    \]
    代入上述 $\xi$ 即得结论.
\end{solution}

%例题4.2.16
\begin{example}
    设 $f(x)$ 在 $(a,b)$ 上三次可导,则存在 $\xi\in(a,b)$,使得
    \[
        f(a)-f(b)+\dfrac{1}{2}(b-a)[f'(a)+f'(b)]
        =f'''(\xi)\dfrac{(b-a)^3}{12}.
    \]
\end{example}
\exampleblank

\begin{solution}
    对函数 $g=f'$ 使用梯形公式的中值型余项.由于
    \[
        \int_a^b g(x)\,\mathrm{d}x=f(b)-f(a),
    \]
    存在 $\xi\in(a,b)$ 使
    \[
        \dfrac{b-a}{2}[g(a)+g(b)]
        -\int_a^b g(x)\,\mathrm{d}x
        =\dfrac{(b-a)^3}{12}g''(\xi).
    \]
    代入 $g=f'$,即得
    \[
        f(a)-f(b)+\dfrac{b-a}{2}[f'(a)+f'(b)]
        =\dfrac{(b-a)^3}{12}f'''(\xi).
    \]
\end{solution}



\subsection{利用几何意义(中值问题)}

Lagrange 中值定理的几何意义即,曲线上任意两点的弦,必与两点之间某点的切线平行.

%例题4.2.17
\begin{example}
    设 $f(x)$ 是可微函数,导函数 $f'(x)$ 严格单调递增,若 $f(a)=f(b)\ (a<b)$,试证明:对一切 $x\in(a,b)$,有
    \[
        f(x)<f(a)=f(b).
    \]
\end{example}
\exampleblank

\begin{solution}
    对任意 $x\in(a,b)$,由 $f'$ 严格递增可知 $f$ 严格凸,因而其割线斜率满足
    \[
        \dfrac{f(x)-f(a)}{x-a}
        <\dfrac{f(b)-f(x)}{b-x}.
    \]
    代入 $f(a)=f(b)$ 并整理,得到 $f(x)<f(a)=f(b)$.
\end{solution}

%例题4.2.18
\begin{example}
    设 $f(x)$ 在区间 $\lbrack 0,1\rbrack$ 上可微,$f(0)=0,f(1)=1$,$k_1,k_2,\ldots,k_n$ 为 $n$ 个正数,证明在区间 $\lbrack 0,1\rbrack$ 上存在一组互不相等的数 $x_1,x_2,\ldots,x_n$,使得
    \[
        \sum\limits_{i=1}^n\dfrac{k_i}{f'(x_i)}
        =\sum\limits_{i=1}^n k_i.
    \]
\end{example}
\exampleblank

\begin{solution}
    若 $f(x)=x$,任取 $n$ 个互不相同的内点即可.以下设 $f$ 不是该线性函数.由
    \[
        f(1)-f(0)=1
    \]
    及中值定理,$f'$ 取到数值 $1$.而且 $f'$ 必须在 $1$ 的两侧都取值;否则 $f(x)-x$ 单调且两端同值,只能恒为零.由导函数的 Darboux 性质,存在 $\delta>0$,使 $(1-\delta,1+\delta)$ 中每个数都是 $f'$ 的值.

    可在该区间内选取互不相同的正数 $r_1,\ldots,r_n$,使
    \[
        \sum\limits_{i=1}^{n}\dfrac{k_i}{r_i}
        =\sum\limits_{i=1}^{n}k_i.
    \]
    例如先在 $1$ 附近任取互不相同的 $r_1,ldots,r_{n-1}$,再由连续性调节 $r_n$;取值充分接近 $1$ 时仍落在上述区间内.分别取 $x_i$ 使 $f'(x_i)=r_i$.因 $r_i$ 互不相同,$x_i$ 也互不相同,代入即得所证等式.
\end{solution}

%例题4.2.19
\begin{example}
    设函数 $f(x)$ 在闭区间 $\lbrack a,b\rbrack$ 上连续,在开区间 $(a,b)$ 内可导,又 $f(x)$ 不是线性函数,且 $f(b)>f(a)$,试证:存在 $\xi\in(a,b)$,使得
    \[
        f'(\xi)>\dfrac{f(b)-f(a)}{b-a}.
    \]
\end{example}
\exampleblank

\begin{solution}
    记割线斜率
    \[
        m=\dfrac{f(b)-f(a)}{b-a}>0,
    \]
    并令 $g(x)=f(x)-f(a)-m(x-a)$.则 $g(a)=g(b)=0$,且 $g$ 不恒为零.取一点 $x_0\in(a,b)$ 使 $g(x_0)\neq0$.

    若 $g(x_0)>0$,对 $g$ 在 $\lbrack a,x_0\rbrack$ 上使用中值定理,得到某点 $\xi$ 满足 $g'(\xi)>0$;若 $g(x_0)<0$,则在 $\lbrack x_0,b\rbrack$ 上使用中值定理,同样得到 $g'(\xi)>0$.于是
    \[
        f'(\xi)=m+g'(\xi)>m.
    \]
\end{solution}

%例题4.2.20
\begin{example}
    设 $f(x)$ 在 $\lbrack a,b\rbrack$ 上连续,在 $(a,b)$ 内可导,$f(x)$ 既不是常数也不是一次函数,试证:存在 $\xi\in(a,b)$,使得
    \[
        |f'(\xi)|>\left|\dfrac{f(b)-f(a)}{b-a}\right|.
    \]
\end{example}
\exampleblank

\begin{solution}
    若 $f(a)=f(b)$,因 $f$ 不是常数,中值定理给出某点 $\xi$ 使 $f'(\xi)\neq0$,结论成立.

    若 $f(b)>f(a)$,上一题直接给出
    \[
        f'(\xi)>\dfrac{f(b)-f(a)}{b-a}>0.
    \]
    若 $f(b)<f(a)$,则对 $-f$ 应用上一题.两种情形都得到题设的严格不等式.
\end{solution}



\subsection{多中值问题}

对于多个中值点的问题显然前面的常规方法失效了,此时要用到的是 Cauchy 中值定理:
\[
    F(b)-F(a)=\dfrac{F'(\xi)}{G'(\xi)}[G(b)-G(a)].
\]
若选取不同的函数 $G$,便可把 $F(b)-F(a)$ 表示成不同的形式.如另取 $G_1(x)$,则存在 $\xi,\eta\in(a,b)$,使得
\[
    \dfrac{F'(\xi)}{G'(\xi)}[G(b)-G(a)]
    =\dfrac{F'(\eta)}{G_1'(\eta)}[G_1(b)-G_1(a)].
\]
一般地,当 $G(x)$ 取 $n$ 个不同的函数时,便可得到含 $n$ 个中值的 $n-1$ 个等式.

%例题4.2.21
\begin{example}
    设 $f(x)$ 在 $\lbrack a,b\rbrack$ 上连续,在 $(a,b)$ 内可导,$f(a)\neq f(b)$,证明:存在 $\xi,\eta\in(a,b)$ 使得
    \[
        f'(\xi)=\dfrac{a+b}{2\eta}f'(\eta).
    \]
\end{example}
\exampleblank

\begin{solution}
    当前题面对端点没有任何限制,结论并不总成立.取
    \[
        a=-1,
        \qquad b=1,
        \qquad f(x)=x.
    \]
    则 $f(a)\neq f(b)$ 且 $f'(x)=1$,但题设等式右端恒为
    \[
        \dfrac{a+b}{2\eta}f'(\eta)=0,
    \]
    不可能等于 $f'(\xi)=1$.若补充 $a+b\neq0$,可分别对 $f$ 与 $x$,$x^2$ 使用 Cauchy 中值定理得到结论;原题应核对这一遗漏条件.
\end{solution}

%例题4.2.22
\begin{example}
    设 $f(x)$ 在 $\lbrack a,b\rbrack\ (b>a>0)$ 上连续,在 $(a,b)$ 内可导,证明:在 $(a,b)$ 内存在两点 $\xi,\eta$,使得
    \[
        f'(\xi)=\dfrac{\eta^2}{ab}f'(\eta).
    \]
\end{example}
\exampleblank

\begin{solution}
    分别对 $f$ 与 $x$,$\dfrac{1}{x}$ 使用 Cauchy 中值定理.存在 $\xi,\eta\in(a,b)$ 使
    \[
        \dfrac{f(b)-f(a)}{b-a}=f'(\xi)
    \]
    以及
    \[
        \dfrac{f(b)-f(a)}{\frac{1}{b}-\frac{1}{a}}
        =\dfrac{f'(\eta)}{-\frac{1}{\eta^2}}.
    \]
    由 $\dfrac{1}{b}-\dfrac{1}{a}=-\dfrac{b-a}{ab}$,第二式化为
    \[
        \dfrac{f(b)-f(a)}{b-a}
        =\dfrac{\eta^2}{ab}f'(\eta).
    \]
    与第一式比较即得结论.
\end{solution}

%例题4.2.23
\begin{example}
    设 $f(x)$ 在 $\lbrack a,b\rbrack$ 上连续,在 $(a,b)$ 内可导,$0<a<b$,证明:在 $(a,b)$ 内存在 $x_1,x_2,x_3$,使得
    \[
        \dfrac{f'(x_1)}{2x_1}
        =(b^2+a^2)\dfrac{f'(x_2)}{4x_3^2}
        =\dfrac{\ln\dfrac{b}{a}}{b^2-a^2}x_3f'(x_3).
    \]
\end{example}
\exampleblank

\begin{solution}
    题面中间一项的分母疑有排印错误.分别对 $f$ 与 $x^2$,$x^4$,$\ln x$ 使用 Cauchy 中值定理,可得 $x_1,x_2,x_3\in(a,b)$,使
    \[
        \dfrac{f(b)-f(a)}{b^2-a^2}
        =\dfrac{f'(x_1)}{2x_1},
    \]
    \[
        \dfrac{f(b)-f(a)}{b^2-a^2}
        =(b^2+a^2)\dfrac{f'(x_2)}{4x_2^3},
    \]
    以及
    \[
        \dfrac{f(b)-f(a)}{b^2-a^2}
        =\dfrac{\ln\frac{b}{a}}{b^2-a^2}x_3f'(x_3).
    \]
    因此能够严格推出的中间项应为 $\dfrac{(b^2+a^2)f'(x_2)}{4x_2^3}$,而不是当前题面中的 $\dfrac{(b^2+a^2)f'(x_2)}{4x_3^2}$.该处需要对照原始资料核对.
\end{solution}

%例题4.2.24
\begin{example}
    设 $f(x)$ 在区间 $\lbrack a,b\rbrack$ 上连续,在 $(a,b)$ 内可导,且 $a\geq0$(或 $b\leq0$),试证:
    \begin{enumerate}
        \item 存在 $x_1,x_2,x_3\in(a,b)$,使得
              \[
                  f'(x_1)
                  =(b+a)\dfrac{f'(x_2)}{2x_2}
                  =(b^2+ba+a^2)\dfrac{f'(x_3)}{3x_3^2};
              \]
        \item 对任意 $n\in\{1,2,\ldots\}$,存在 $x_1,x_2,\ldots,x_n\in(a,b)$,使得
              \[
                  \begin{aligned}
                      f'(x_1)
                       & =(b+a)\dfrac{f'(x_2)}{2x_2}
                      =\cdots                                        \\
                       & =(b^{n-1}+b^{n-2}a+\cdots+ba^{n-2}+a^{n-1})
                      \dfrac{f'(x_n)}{nx_n^{n-1}}.
                  \end{aligned}
              \]
    \end{enumerate}
\end{example}
\exampleblank

\begin{solution}
    对每个 $j=1,2,\ldots,n$,对 $f$ 与 $x^j$ 使用 Cauchy 中值定理,存在 $x_j\in(a,b)$ 使
    \[
        \dfrac{f(b)-f(a)}{b^j-a^j}
        =\dfrac{f'(x_j)}{jx_j^{j-1}}.
    \]
    又
    \[
        b^j-a^j=(b-a)(b^{j-1}+b^{j-2}a+\cdots+a^{j-1}),
    \]
    故
    \[
        \dfrac{f(b)-f(a)}{b-a}
        =(b^{j-1}+b^{j-2}a+\cdots+a^{j-1})
        \dfrac{f'(x_j)}{jx_j^{j-1}}.
    \]
    依次取 $j=1,2,3$ 得第 (1) 问,取 $j=1,2,\ldots,n$ 得第 (2) 问.
\end{solution}

%例题4.2.25
\begin{example}
    设 $f\in C(\lbrack a,b\rbrack)$,且在 $(a,b)$ 上可导,若 $f(a)=f(b)=1$,则存在 $\xi,\eta\in(a,b)$,使得
    \[
        e^{\eta-\xi}[f(\eta)+f'(\eta)]=1.
    \]
\end{example}
\exampleblank

\begin{solution}
    对函数 $F(x)=e^x f(x)$ 与 $G(x)=e^x$ 使用 Cauchy 中值定理.由于 $f(a)=f(b)=1$,存在 $\eta\in(a,b)$ 使
    \[
        \dfrac{F(b)-F(a)}{G(b)-G(a)}
        =\dfrac{F'(\eta)}{G'(\eta)}
        =f(\eta)+f'(\eta).
    \]
    左端等于 $1$,故 $f(\eta)+f'(\eta)=1$.取 $\xi=\eta$,便有
    \[
        e^{\eta-\xi}[f(\eta)+f'(\eta)]=1.
    \]
\end{solution}

%例题4.2.26
\begin{example}
    设 $f\in C(\lbrack a,b\rbrack)$,且在 $(a,b)$ 上可导,若 $f(0)=0,f(1)=1$,则存在 $\xi,\eta\in(a,b)$,使得
    \[
        f'(\xi)f'(\eta)=1,
        \qquad \xi\neq\eta.
    \]
\end{example}
\exampleblank

\begin{solution}
    题设中的定义区间应至少包含 $\lbrack0,1\rbrack$.若 $f(x)=x$,任取两个不同的内点即可.否则,由中值定理及导函数的 Darboux 性质,$f'$ 在数值 $1$ 的两侧取值,并包含某个区间 $(1-\delta,1+\delta)$.

    取 $r>1$ 充分接近 $1$,使 $r$ 与 $\dfrac{1}{r}$ 都在该区间内.于是存在 $\xi,\eta\in(0,1)$ 使
    \[
        f'(\xi)=r,
        \qquad
        f'(\eta)=\dfrac{1}{r}.
    \]
    两导数值不同,故 $\xi\neq\eta$,并且 $f'(\xi)f'(\eta)=1$.
\end{solution}



\subsection{函数与导函数的关系}

设 $f(x)$ 在 $(a,+\infty)$ 上可导,如果只有极限 $\displaystyle\lim\limits_{x\to+\infty}f'(x)=B$ 存在,那么当 $x\to+\infty$ 时,$f(x)$ 不一定存在极限,例如 $f(x)=\sin(\ln x)$;反之,即使存在 $\displaystyle\lim\limits_{x\to+\infty}f(x)$,也不一定存在 $\displaystyle\lim\limits_{x\to+\infty}f'(x)$,例如 $f(x)=\dfrac{\sin x^2}{x}$.但若二者都存在,那么必有 $\displaystyle\lim\limits_{x\to+\infty}f'(x)=0$.

%例题4.2.27
\begin{example}
    设 $f(x)$ 在 $(a,+\infty)$ 上可导,若 $\displaystyle\lim\limits_{x\to+\infty}f'(x)=B>0$,则:
    \begin{enumerate}
        \item $\displaystyle\lim\limits_{x\to+\infty}f(x)=+\infty$;
        \item $f(x)\sim Bx\ (x\to+\infty)$.
    \end{enumerate}
\end{example}
\exampleblank

\begin{solution}
    取 $X>a$,使得当 $x>X$ 时 $f'(x)>\dfrac{B}{2}$.于是由中值定理,
    \[
        f(x)-f(X)>\dfrac{B}{2}(x-X),
    \]
    故 $f(x)\to+\infty$.

    又因 $x\to+\infty$ 时 $f(x)$ 与 $x$ 都趋于 $+\infty$,由 Cauchy 中值定理或 L'Hospital 法则,
    \[
        \lim\limits_{x\to+\infty}\dfrac{f(x)}{x}
        =\lim\limits_{x\to+\infty}f'(x)=B.
    \]
    因此 $f(x)\sim Bx$.
\end{solution}

%例题4.2.28
\begin{example}[\markstar]
    设 $f\in C^{(1)}((a,+\infty))$,若 $f'(x)$ 在 $(a,+\infty)$ 上一致连续,且存在 $\displaystyle\lim\limits_{x\to+\infty}f(x)$,则
    \[
        \lim\limits_{x\to+\infty}f'(x)=0.
    \]
\end{example}
\exampleblank

\begin{solution}
    反设结论不成立,则存在 $\varepsilon>0$ 与数列 $x_n\to+\infty$,使 $|f'(x_n)|\geq\varepsilon$.由 $f'$ 一致连续,存在与 $n$ 无关的 $\delta>0$,使
    \[
        |x-x_n|<\delta
        \quad\Longrightarrow\quad
        |f'(x)-f'(x_n)|<\dfrac{\varepsilon}{2}.
    \]
    取子列可设 $f'(x_n)\geq\varepsilon$,或恒有 $f'(x_n)\leq-\varepsilon$.在相应的长度为 $\delta$ 的区间上,$f'$ 保持同号且绝对值不小于 $\dfrac{\varepsilon}{2}$,于是 $f$ 的函数值变化至少为 $\dfrac{\varepsilon\delta}{2}$.这与 $f(x)$ 在无穷远处有有限极限,从而满足 Cauchy 条件矛盾.故 $f'(x)\to0$.
\end{solution}

%例题4.2.29
\begin{example}[\markstar]
    设 $f(x)$ 在 $\lbrack 0,+\infty)$ 上三阶可导,且有
    \[
        \lim\limits_{x\to+\infty}f(x)=A,
        \qquad
        \lim\limits_{x\to+\infty}f'''(x)=0,
    \]
    则
    \[
        \lim\limits_{x\to+\infty}f'(x)
        =\lim\limits_{x\to+\infty}f''(x)=0.
    \]
\end{example}
\exampleblank

\begin{solution}
    先证 $f''(x)\to0$.固定 $h>0$.对 $f(x+h)$ 与 $f(x+2h)$ 在 $x$ 处作三阶 Taylor 展开并使用 Lagrange 余项,消去一次项可得
    \[
        f(x+2h)-2f(x+h)+f(x)
        =h^2f''(x)+R_x,
    \]
    其中当 $x\to+\infty$ 时,$\dfrac{|R_x|}{h^2}\to0$,因为展开区间上的 $f'''$ 一致趋于零.左端也因 $f(x)\to A$ 而趋于零,故 $f''(x)\to0$.

    再由二阶 Taylor 公式,存在 $\theta_x\in(0,1)$ 使
    \[
        f(x+h)-f(x)
        =hf'(x)+\dfrac{h^2}{2}f''(x+\theta_xh).
    \]
    令 $x\to+\infty$,左端与最后一项都趋于零,故 $f'(x)\to0$.
\end{solution}

%例题4.2.30
\begin{example}[\markstar]
    试证明下列命题:
    \begin{enumerate}
        \item 设 $f(x)$ 在 $(0,+\infty)$ 上可导,$a>0$,若
              \[
                  \lim\limits_{x\to+\infty}[af(x)+f'(x)]=l,
              \]
              则
              \[
                  \lim\limits_{x\to+\infty}f(x)=\dfrac{l}{a};
              \]
        \item 设 $f(x)$ 在 $(0,+\infty)$ 上可导,$a>0$,若有
              \[
                  \lim\limits_{x\to+\infty}[af(x)+2\sqrt{x}f'(x)]=l,
              \]
              则
              \[
                  \lim\limits_{x\to+\infty}f(x)=\dfrac{l}{a};
              \]
        \item 设 $f(x)$ 在 $(0,+\infty)$ 上二次可导,若有
              \[
                  \lim\limits_{x\to+\infty}[f(x)+f'(x)+f''(x)]=l,
              \]
              则
              \[
                  \lim\limits_{x\to+\infty}f(x)=l.
              \]
    \end{enumerate}
\end{example}
\exampleblank

\begin{solution}
    记趋于零的余项为 $r(x)$.

    \begin{enumerate}
        \item 令 $u(x)=f(x)-\dfrac{l}{a}$,则
              \[
                  u'(x)+au(x)=r(x),
                  \qquad r(x)\to0.
              \]
              乘以积分因子 $e^{ax}$ 并从固定点 $x_0$ 积分,得到
              \[
                  u(x)=e^{-a(x-x_0)}u(x_0)
                  +\int_{x_0}^{x}e^{-a(x-t)}r(t)\,\mathrm{d}t.
              \]
              第一项趋于零;将积分按 $r(t)$ 在有限区间和尾部拆分,也可知第二项趋于零.因此 $f(x)\to\dfrac{l}{a}$.

        \item 令 $F(t)=f(t^2)$.则
              \[
                  F'(t)=2t f'(t^2),
              \]
              因而 $aF(t)+F'(t)\to l$.由第 (1) 问,$F(t)\to\dfrac{l}{a}$,即 $f(x)\to\dfrac{l}{a}$.

        \item 令 $u=f-l$,则
              \[
                  u''+u'+u=r(x),
                  \qquad r(x)\to0.
              \]
              该稳定线性方程的常数变易公式为
              \[
                  \begin{aligned}
                      u(x)={} & e^{-\frac{x-x_0}{2}}
                      \left[C_1\cos\dfrac{\sqrt3(x-x_0)}{2}
                      +C_2\sin\dfrac{\sqrt3(x-x_0)}{2}\right]                      \\
                              & +\dfrac{2}{\sqrt3}\int_{x_0}^{x}e^{-\frac{x-t}{2}}
                      \sin\dfrac{\sqrt3(x-t)}{2}\,r(t)\,\mathrm{d}t.
                  \end{aligned}
              \]
              齐次项趋于零;卷积项同样由 $r(t)\to0$ 趋于零,故 $u(x)\to0$,即 $f(x)\to l$.
    \end{enumerate}
\end{solution}

%例题4.2.31
\begin{example}
    设 $f(x)$ 在 $(0,+\infty)$ 上可导,若 $\displaystyle\lim\limits_{x\to+\infty}f'(x)=l$,则对任意的 $A$ 值,有
    \[
        \lim\limits_{x\to+\infty}[f(x+A)-f(x)]=lA.
    \]
\end{example}
\exampleblank

\begin{solution}
    若 $A=0$,结论显然.若 $A>0$,由 Lagrange 中值定理,存在 $\xi_x\in(x,x+A)$ 使
    \[
        f(x+A)-f(x)=Af'(\xi_x).
    \]
    若 $A<0$,同样存在 $\xi_x\in(x+A,x)$ 使该式成立.两种情形下都有 $\xi_x\to+\infty$,故令 $x\to+\infty$ 得
    \[
        f(x+A)-f(x)\longrightarrow Al.
    \]
\end{solution}

%例题4.2.32
\begin{example}[\markstar]
    设 $f(x)$ 在 $\lbrack 0,+\infty)$ 上可微,$f(0)=0$,且
    \[
        |f'(x)|\leq A|f(x)|,
        \qquad x\in\lbrack 0,+\infty),
    \]
    则 $f\equiv0$.
\end{example}
\exampleblank

\begin{solution}
    由 $f(0)=0$,对任意 $x\geq0$ 有
    \[
        |f(x)|\leq\int_0^x|f'(t)|\,\mathrm{d}t
        \leq A\int_0^x|f(t)|\,\mathrm{d}t.
    \]
    令 $G(x)=\displaystyle\int_0^x|f(t)|\,\mathrm{d}t$,则 $G(0)=0$ 且 $G'(x)\leq AG(x)$.于是
    \[
        \left(e^{-Ax}G(x)\right)'\leq0.
    \]
    但 $e^{-Ax}G(x)\geq0$ 且其在 $0$ 处为零,故 $G\equiv0$,进而 $f\equiv0$.
\end{solution}

%例题4.2.33
\begin{example}
    设 $f(x),g(x)$ 在 $\lbrack a,b\rbrack$ 上连续,$g(x)$ 在 $(a,b)$ 内可微,$g(a)=0$,若有实数 $\lambda\neq0$,使得
    \[
        |g(x)f(x)+\lambda g'(x)|\leq|g(x)|,
        \qquad x\in(a,b),
    \]
    成立,则 $g\equiv0$.
\end{example}
\exampleblank

\begin{solution}
    由 $f$ 在紧区间上的连续性,存在 $M>0$ 使 $|f(x)|\leq M$.题设不等式给出
    \[
        |\lambda|\,|g'(x)|
        \leq |g(x)|+|g(x)f(x)|
        \leq(M+1)|g(x)|.
    \]
    因此
    \[
        |g'(x)|\leq C|g(x)|,
        \qquad C=\dfrac{M+1}{|\lambda|}.
    \]
    结合 $g(a)=0$,应用上一题的结论便得 $g\equiv0$.
\end{solution}

%例题4.2.34
\begin{example}
    证明:在 $\mathbb{R}$ 上的任何可微函数都不可能满足下列函数方程:
    \begin{enumerate}
        \item $g(g(x))=-x^3+x+1$;
        \item $g(g(x))=-x^3+x^2\pm1$;
        \item $g(g(x))=x^2-3x+3$.
    \end{enumerate}
\end{example}
\exampleblank

\begin{solution}
    记右端多项式为 $P(x)$.若 $P=g\circ g$,则 $P$ 的任一不动点 $x_0$ 都满足 $g(x_0)$ 也是 $P$ 的不动点,并且
    \[
        P'(x_0)=g'(g(x_0))g'(x_0).
    \]

    \begin{enumerate}
        \item $P(x)=-x^3+x+1$ 只有不动点 $x=1$,故 $g(1)=1$.但
              \[
                  P'(1)=[g'(1)]^2\geq0,
              \]
              而直接计算得 $P'(1)=-2$,矛盾.

        \item 对 $P(x)=-x^3+x^2+1$,唯一不动点为 $1$,且 $P'(1)=-1$,同样矛盾.对 $P(x)=-x^3+x^2-1$,不动点方程
              \[
                  x^3-x^2+x+1=0
              \]
              的左端严格递增,故只有一个实根 $r\in(-1,0)$.于是 $g(r)=r$,但
              \[
                  P'(r)=-3r^2+2r<0,
              \]
              不可能等于 $[g'(r)]^2$.

        \item $P(x)=x^2-3x+3$ 的不动点为 $1,3$.若 $g$ 固定其中某一点,则该点处 $P'$ 应为平方而非负;但 $P'(1)=-1$,所以只能有 $g(1)=3$,$g(3)=1$.链式法则又给出
              \[
                  P'(1)=g'(3)g'(1)=P'(3),
              \]
              这与 $P'(1)=-1$,$P'(3)=3$ 矛盾.
    \end{enumerate}
\end{solution}


\newpage

\section{Taylor 公式}


\subsection{中值等式}

%例题4.3.1
\begin{example}
    设 $f(x)$ 在 $\lbrack a,b\rbrack$ 上三次可导,试证:存在 $c\in(a,b)$,使得
    \[
        f(b)=f(a)+f'\left(\dfrac{a+b}{2}\right)(b-a)
        +\dfrac{1}{24}f'''(c)(b-a)^3.
    \]
\end{example}
\exampleblank

\begin{solution}
    记 $m=\dfrac{a+b}{2}$.这是中心 Taylor 公式的中值型余项.具体地,对 $f$ 在 $m$ 两侧的差构造三次 Hermite 插值余项,连续使用 Rolle 定理三次,得到某个 $c\in(a,b)$ 使
    \[
        f(b)-f(a)
        =(b-a)f'(m)+\dfrac{(b-a)^3}{24}f'''(c).
    \]
    移项即为题设公式.
\end{solution}

%例题4.3.2
\begin{example}
    设 $f\in C^{(3)}(\lbrack -1,1\rbrack)$,且 $f(-1)=0,f(1)=1,f'(0)=0$,则存在 $\xi\in(-1,1)$,使得
    \[
        f'''(\xi)=3.
    \]
\end{example}
\exampleblank

\begin{solution}
    作三次多项式 $P$,使
    \[
        P(-1)=f(-1),
        \quad P(0)=f(0),
        \quad P'(0)=f'(0),
        \quad P(1)=f(1).
    \]
    写成 $P(x)=Ax^3+Bx^2+Cx+D$.由 $C=f'(0)=0$ 以及
    \[
        P(1)-P(-1)=f(1)-f(-1)=1
    \]
    可得 $2A=1$,所以 $P'''(x)=6A=3$.

    函数 $f-P$ 在 $-1,0,1$ 为零,且在 $0$ 处导数也为零.由带重节点的 Rolle 定理,存在 $\xi\in(-1,1)$ 使
    \[
        (f-P)'''(\xi)=0.
    \]
    故 $f'''(\xi)=P'''(\xi)=3$.
\end{solution}

%例题4.3.3
\begin{example}
    设 $f(x)$ 在 $\lbrack -1,1\rbrack$ 上二次可导,若有 $f(-1)=0,f(0)=0,f(1)=1$,则存在 $\xi\in(-1,1)$,使得
    \[
        f''(\xi)=1.
    \]
\end{example}
\exampleblank

\begin{solution}
    通过三点 $(-1,0)$,$(0,0)$,$(1,1)$ 的二次插值多项式为
    \[
        P(x)=\dfrac{x^2+x}{2},
    \]
    故 $P''(x)=1$.函数 $f-P$ 在 $-1,0,1$ 三点为零,连续使用两次 Rolle 定理,存在 $\xi\in(-1,1)$ 使 $(f-P)''(\xi)=0$.因此 $f''(\xi)=1$.
\end{solution}

%例题4.3.4
\begin{example}
    设 $f(x)$ 在 $\lbrack -1,1\rbrack$ 上三次可导,若有 $f(-1)=f(0)=0,f(1)=1,f'(0)=0$,则存在 $\xi\in(-1,1)$,使得
    \[
        f'''(\xi)\geq 3.
    \]
\end{example}
\exampleblank

\begin{solution}
    上一题组中的三次 Hermite 插值论证在这里同样适用.事实上,$f(-1)=0$,$f(1)=1$,$f'(0)=0$ 已足以推出存在 $\xi\in(-1,1)$ 使
    \[
        f'''(\xi)=3.
    \]
    因此当然有 $f'''(\xi)\geq3$;附加条件 $f(0)=0$ 不影响结论.
\end{solution}



\subsection{界的估计}

%例题4.3.5
\begin{example}
    设 $f(x)$ 二次可微,且 $f(0)=f(1)=0$,又
    $\displaystyle\max\limits_{0\leq x\leq1}f(x)=2$,试证:
    \[
        \min\limits_{0\leq x\leq1}f''(x)\leq-16.
    \]
\end{example}
\exampleblank

\begin{solution}
    设 $f$ 在 $x_0\in(0,1)$ 处取得最大值 $2$,则 $f'(x_0)=0$.分别在区间 $\lbrack0,x_0\rbrack$ 与 $\lbrack x_0,1\rbrack$ 使用二阶 Taylor 公式,存在 $\xi_1\in(0,x_0)$,$\xi_2\in(x_0,1)$ 使
    \[
        f''(\xi_1)=-\dfrac{4}{x_0^2},
        \qquad
        f''(\xi_2)=-\dfrac{4}{(1-x_0)^2}.
    \]
    由于 $x_0$ 与 $1-x_0$ 至少有一个不超过 $\dfrac{1}{2}$,上述两数至少有一个不大于 $-16$.故
    \[
        \min\limits_{0\leq x\leq1}f''(x)\leq-16.
    \]
\end{solution}

%例题4.3.6
\begin{example}
    已知 $f(x)$ 在 $\lbrack 0,1\rbrack$ 上二阶可导,且 $f(0)=0,f(1)=3$,又
    $\displaystyle\min\limits_{x\in\lbrack 0,1\rbrack}f(x)=-1$,证明:存在 $c\in(0,1)$,使得
    \[
        f''(c)\geq18.
    \]
\end{example}
\exampleblank

\begin{solution}
    设 $f$ 在 $x_0\in(0,1)$ 处取得最小值 $-1$,则 $f'(x_0)=0$.由二阶 Taylor 公式,存在 $\xi_1\in(0,x_0)$,$\xi_2\in(x_0,1)$ 使
    \[
        f''(\xi_1)=\dfrac{2}{x_0^2},
        \qquad
        f''(\xi_2)=\dfrac{8}{(1-x_0)^2}.
    \]
    而对 $0<x_0<1$,
    \[
        \max\left\{\dfrac{2}{x_0^2},\dfrac{8}{(1-x_0)^2}\right\}\geq18;
    \]
    等号在 $x_0=\dfrac{1}{3}$ 时取得.因此至少有一点 $c\in(0,1)$ 满足 $f''(c)\geq18$.
\end{solution}

%例题4.3.7
\begin{example}[\markstar]
    设 $f(x)$ 在 $\lbrack 0,1\rbrack$ 上二次可导,若有
    \[
        |f(x)|\leq A,\qquad |f''(x)|\leq B,\qquad x\in\lbrack 0,1\rbrack,
    \]
    则
    \[
        |f'(x)|\leq2A+\dfrac{B}{2}.
    \]
\end{example}
\exampleblank

\begin{solution}
    在点 $x\in\lbrack0,1\rbrack$ 分别向端点 $0,1$ 使用二阶 Taylor 公式,存在相应的中间点 $\xi_0,\xi_1$ 使
    \[
        f(1)-f(0)
        =f'(x)+\dfrac{(1-x)^2}{2}f''(\xi_1)
        -\dfrac{x^2}{2}f''(\xi_0).
    \]
    因此
    \[
        \begin{aligned}
            |f'(x)|
             & \leq |f(1)-f(0)|
            +\dfrac{B}{2}\left[(1-x)^2+x^2\right] \\
             & \leq2A+\dfrac{B}{2}.
        \end{aligned}
    \]
\end{solution}

%例题4.3.8
\begin{example}
    设 $f(x)$ 在 $\lbrack 0,1\rbrack$ 上二阶可导,且有 $f(0)=f(1)$,$|f''(x)|\leq M$,则
    \[
        |f'(x)|\leq\dfrac{M}{2}.
    \]
\end{example}
\exampleblank

\begin{solution}
    使用上一题中的恒等式,并代入 $f(0)=f(1)$,得到
    \[
        |f'(x)|
        \leq\dfrac{M}{2}\left[x^2+(1-x)^2\right]
        \leq\dfrac{M}{2}.
    \]
\end{solution}

%例题4.3.9
\begin{example}[\markstar]
    设 $f(x)$ 在 $(-\infty,+\infty)$ 上二次可导,且有
    \[
        M_k=\sup\{|f^{(k)}(x)|:-\infty<x<+\infty\}<+\infty.
    \]
    则
    \[
        M_1\leq\sqrt{2M_0M_2}.
    \]
\end{example}
\exampleblank

\begin{solution}
    对任意 $h>0$,在 $x$ 的两侧使用二阶 Taylor 公式并相减,得到
    \[
        |f'(x)|\leq\dfrac{M_0}{h}+\dfrac{M_2h}{2}.
    \]
    若 $M_2>0$,取 $h=\sqrt{\dfrac{2M_0}{M_2}}$,便有
    \[
        |f'(x)|\leq\sqrt{2M_0M_2}.
    \]
    再对 $x$ 取上确界即得结论.$M_0=0$ 或 $M_2=0$ 时可直接由函数为常数或令 $h\to+\infty$ 得到同一结论.
\end{solution}

%例题4.3.10
\begin{example}[\markstar]
    设 $f(x)$ 在 $(-\infty,+\infty)$ 上 $m$ 次可导,且有
    \[
        M_k=\sup\{|f^{(k)}(x)|:-\infty<x<+\infty\}<+\infty.
    \]
    则
    \[
        M_k\leq2^{\frac{k(m-k)}{2}}M_0^{1-\frac{k}{m}}M_m^{\frac{k}{m}},
        \qquad k=1,2,\ldots,m-1.
    \]
\end{example}
\exampleblank

\begin{solution}
    将上一题应用于 $f^{(j-1)}$,可得
    \[
        M_j^2\leq2M_{j-1}M_{j+1},
        \qquad j=1,2,\ldots,m-1.
    \]
    当各 $M_j>0$ 时,令 $A_j=\ln M_j$.上式等价于
    \[
        2A_j\leq A_{j-1}+A_{j+1}+\ln2.
    \]
    离散凸性比较给出
    \[
        A_k\leq\left(1-\dfrac{k}{m}\right)A_0
        +\dfrac{k}{m}A_m+\dfrac{k(m-k)}{2}\ln2.
    \]
    两端取指数,便得到
    \[
        M_k\leq2^{\frac{k(m-k)}{2}}
        M_0^{1-\frac{k}{m}}M_m^{\frac{k}{m}}.
    \]
    若某个 $M_j=0$,用极限或直接由相应导数恒为零处理即可.
\end{solution}



\subsection{中值点的极限}

%例题4.3.11
\begin{example}
    证明:对每个 $x>0$,均存在 $\theta(x)\in\lbrack \dfrac{1}{4},\dfrac{1}{2}\rbrack$,使得
    \[
        \sqrt{x+1}-\sqrt{x}=\dfrac{1}{2\sqrt{x+\theta(x)}}.
    \]
    且
    \[
        \lim\limits_{x\to0^+}\theta(x)=\dfrac{1}{4},
        \qquad
        \lim\limits_{x\to+\infty}\theta(x)=\dfrac{1}{2}.
    \]
\end{example}
\exampleblank

\begin{solution}
    由有理化公式,
    \[
        \sqrt{x+1}-\sqrt{x}
        =\dfrac{1}{\sqrt{x+1}+\sqrt{x}}.
    \]
    因此只需令
    \[
        2\sqrt{x+\theta(x)}=\sqrt{x+1}+\sqrt{x},
    \]
    即
    \[
        \theta(x)=\dfrac{1}{4}
        +\dfrac{\sqrt{x^2+x}-x}{2}.
    \]
    因 $0<\sqrt{x^2+x}-x<\dfrac{1}{2}$,故 $\dfrac{1}{4}<\theta(x)<\dfrac{1}{2}$.同时
    \[
        \sqrt{x^2+x}-x
        =\dfrac{x}{\sqrt{x^2+x}+x},
    \]
    从而当 $x\to0^+$ 时该差趋于 $0$,当 $x\to+\infty$ 时趋于 $\dfrac{1}{2}$.于是所述两个极限分别为 $\dfrac{1}{4}$ 与 $\dfrac{1}{2}$.
\end{solution}

%例题4.3.12
\begin{example}
    设 $f(x)$ 在点 $a$ 的某个邻域上 $n+2$ 阶连续可微,且 $f^{(n+2)}(a)\neq0$,并有
    \[
        f(a+h)=f(a)+f'(a)h+\cdots+\dfrac{f^{(n)}(a)}{n!}h^n
        +\dfrac{f^{(n+1)}(a+\theta h)}{(n+1)!}h^{n+1},
    \]
    其中 $h\in(0,\delta)$,$0<\theta<1$.证明:
    \[
        \lim\limits_{h\to0}\theta=\dfrac{1}{n+2}.
    \]
\end{example}
\exampleblank

\begin{solution}
    在 $a$ 处将等式两端展开到 $n+2$ 阶.左端为
    \[
        \begin{aligned}
            f(a+h)={} & f(a)+f'(a)h+\cdots
            +\dfrac{f^{(n+1)}(a)}{(n+1)!}h^{n+1}                         \\
                      & +\dfrac{f^{(n+2)}(a)}{(n+2)!}h^{n+2}+o(h^{n+2}).
        \end{aligned}
    \]
    又因 $0<\theta<1$,一致地有
    \[
        f^{(n+1)}(a+\theta h)
        =f^{(n+1)}(a)+\theta h f^{(n+2)}(a)+o(h).
    \]
    比较 $h^{n+2}$ 项,并利用 $f^{(n+2)}(a)\neq0$,得到
    \[
        \dfrac{\theta}{(n+1)!}\longrightarrow\dfrac{1}{(n+2)!},
    \]
    故 $\theta\to\dfrac{1}{n+2}$.
\end{solution}

%例题4.3.13
\begin{example}
    设 $f(x)$ 在 $(x_0-\delta,x_0+\delta)$ 上 $n$ 阶连续可微,已知
    \[
        f^{(2)}(x_0)=f^{(3)}(x_0)=\cdots=f^{(n-1)}(x_0)=0,
        \qquad f^{(n)}(x_0)\neq0,
    \]
    若 $0<|h|<\delta$ 时有
    \[
        \dfrac{f(x_0+h)-f(x_0)}{h}=f'(x_0+\theta h),
    \]
    证明:
    \[
        \lim\limits_{h\to0}\theta
        =\sqrt[n-1]{\dfrac{1}{n}}.
    \]
\end{example}
\exampleblank

\begin{solution}
    由题设的各阶导数消失条件,
    \[
        \dfrac{f(x_0+h)-f(x_0)}{h}
        =f'(x_0)+\dfrac{f^{(n)}(x_0)}{n!}h^{n-1}+o(h^{n-1}).
    \]
    另一方面,因 $0<\theta<1$,
    \[
        f'(x_0+\theta h)
        =f'(x_0)+\dfrac{f^{(n)}(x_0)}{(n-1)!}
        \theta^{n-1}h^{n-1}+o(h^{n-1}).
    \]
    比较两式并利用 $f^{(n)}(x_0)\neq0$,可得
    \[
        \theta^{n-1}\longrightarrow\dfrac{1}{n}.
    \]
    由于 $0<\theta<1$,故 $\theta\to\sqrt[n-1]{\dfrac{1}{n}}$.
\end{solution}

%例题4.3.14
\begin{example}
    设 $f(x)$ 在点 $a$ 处二阶可导,且 $f''(a)=0$.证明当 $h$ 充分小时,有等式
    \[
        f(a+h)-f(a)=f'(a+\theta h)h,
    \]
    其中 $\theta\in(0,1)$,并且
    \[
        \lim\limits_{h\to0}\theta=\dfrac{1}{2}.
    \]
\end{example}
\exampleblank

\begin{solution}
    本题当前把条件写成了 $f''(a)=0$,在此条件下所述极限一般不成立.取
    \[
        f(x)=x+x^3,
        \qquad a=0.
    \]
    则 $f''(0)=0$.中值等式要求
    \[
        h+h^3=h\left(1+3\theta^2h^2\right),
    \]
    因而对每个 $h\neq0$ 都有 $\theta=\dfrac{1}{\sqrt3}$,并不趋于 $\dfrac{1}{2}$.标准结论通常要求 $f''(a)\neq0$;此时比较二阶 Taylor 展开即可得到 $\theta\to\dfrac{1}{2}$.原题符号需要核对.
\end{solution}



\subsection{极值问题}

设 $f(x)$ 在 $x_0$ 处二阶可导,且 $f'(x_0)=0$.若 $f''(x_0)<0$($>0$),则 $x_0$ 为 $f(x)$ 的极大(小)值点.

%例题4.3.15
\begin{example}
    设 $f\in C(\lbrack a,b\rbrack)$,且 $f\in C^{(2)}((a,b))$,若有
    \[
        f''(x)=e^x f(x),\qquad f(a)=f(b)=0,
    \]
    则 $f(x)\equiv0$.
\end{example}
\exampleblank

\begin{solution}
    反设 $f$ 不恒为零.若 $f$ 在某个内点取得正的最大值,则该点有 $f''\leq0$,但方程给出 $f''=e^x f>0$,矛盾;若 $f$ 取得负的最小值,则该点有 $f''\geq0$,而方程给出 $f''=e^x f<0$,仍矛盾.结合 $f(a)=f(b)=0$,可知 $f$ 既不能为正也不能为负,故 $f\equiv0$.
\end{solution}

%例题4.3.16
\begin{example}
    设 $f\in C^{(2)}((-\infty,+\infty))$,若有
    \[
        [f(0)]^2+[f'(0)]^2=4,
        \qquad |f(x)|=1\quad(-\infty<x<+\infty),
    \]
    则存在 $\xi$,使得
    \[
        f(\xi)+f''(\xi)=0.
    \]
\end{example}
\exampleblank

\begin{solution}
    当前题设彼此矛盾.由于 $f$ 连续且 $|f(x)|=1$ 对所有实数成立,而实轴连通,$f$ 只能恒等于 $1$ 或恒等于 $-1$.于是 $f'(0)=0$,从而
    \[
        [f(0)]^2+[f'(0)]^2=1,
    \]
    不可能等于 $4$.因此原题中的 $|f(x)|=1$ 很可能是遗漏了不等号,需要核对原始资料.
\end{solution}

%例题4.3.17
\begin{example}
    设 $f(x),g(x)$ 定义在 $(-\infty,+\infty)$ 上,且 $f(x)$ 二次可导,又有
    \[
        f(a)=f(b)=0\quad(a<b),
    \]
    且
    \[
        f''(x)+f'(x)g(x)-f(x)=0,
        \qquad x\in(-\infty,+\infty),
    \]
    则 $f(x)\equiv0$.
\end{example}
\exampleblank

\begin{solution}
    当前题面没有要求 $g$ 连续,结论并不成立.固定 $a<b$,令
    \[
        f(x)=
        \begin{cases}
            0,       & x\leq b, \\
            (x-b)^3, & x>b,
        \end{cases}
    \]
    并令 $g(x)=0$($x\leq b$),而当 $x>b$ 时令
    \[
        g(x)=\dfrac{x-b}{3}-\dfrac{2}{x-b}.
    \]
    则 $f$ 二次可导,$f(a)=f(b)=0$,且逐点满足题设微分方程,但 $f$ 并不恒为零.

    若补充 $g$ 连续,则线性方程初值问题具有唯一性.由最大值原理先得 $f\equiv0$ 于 $\lbrack a,b\rbrack$,再利用 $f(b)=f'(b)=0$ 及唯一性即可推出全轴上 $f\equiv0$.因此原题需要补充 $g$ 的正则性条件.
\end{solution}

%例题4.3.18
\begin{example}
    设 $f\in C^{(2)}(\lbrack 0,a\rbrack)$,且 $x_0\in(0,a)$ 是 $f(x)$ 的极小值点,则
    \[
        |f'(0)|+|f'(a)|
        \leq a\max\limits_{x\in\lbrack 0,a\rbrack}|f''(x)|.
    \]
\end{example}
\exampleblank

\begin{solution}
    记
    \[
        M=\max\limits_{x\in\lbrack0,a\rbrack}|f''(x)|.
    \]
    由于 $x_0$ 是内点极小值点,$f'(x_0)=0$.由 Lagrange 中值定理,
    \[
        |f'(0)|=|f'(0)-f'(x_0)|\leq Mx_0,
    \]
    \[
        |f'(a)|=|f'(a)-f'(x_0)|\leq M(a-x_0).
    \]
    两式相加即得题设不等式.
\end{solution}


\newpage

\section{凹凸性与不等式}


\subsection{凹凸性}

\begin{definition}
    设 $f(x)$ 在区间 $I$ 上有定义,$f(x)$ 在 $I$ 上称为凸函数,当且仅当对任意 $x_1,x_2\in I$ 以及任意 $\lambda\in(0,1)$,有
    \[
        f(\lambda x_1+(1-\lambda)x_2)
        \leq\lambda f(x_1)+(1-\lambda)f(x_2).
    \]
\end{definition}

\begin{definition}
    设 $f(x)$ 在区间 $I$ 上有定义,$f(x)$ 在 $I$ 上称为凸函数,当且仅当对任意 $x_1,x_2\in I$,有
    \[
        f\left(\dfrac{x_1+x_2}{2}\right)
        \leq\dfrac{f(x_1)+f(x_2)}{2}.
    \]
\end{definition}

\begin{definition}
    设 $f(x)$ 在区间 $I$ 上有定义,$f(x)$ 在 $I$ 上称为凸函数,当且仅当对任意 $x_1,x_2,\ldots,x_n\in I$,有
    \[
        f\left(\dfrac{x_1+x_2+\cdots+x_n}{n}\right)
        \leq\dfrac{f(x_1)+f(x_2)+\cdots+f(x_n)}{n}.
    \]
\end{definition}

\begin{definition}
    设 $f(x)$ 在区间 $I$ 上有定义,当且仅当曲线 $y=f(x)$ 的切线保持在曲线以下,称 $f(x)$ 为凸函数.
\end{definition}

\begin{lemma}[\markstar]
    设 $f(x)$ 是区间 $I$ 上的凸函数,则对 $I$ 的任一内点 $x$,单侧导数 $f'_+(x),f'_-(x)$ 皆存在,皆为增函数,且
    \[
        f'_-(x)\leq f'_+(x),
        \qquad x\in\mathring I.
    \]
\end{lemma}

因此,更一般地,有如下定理.

\begin{theorem}[\markstar]
    $f$ 为凸函数,当且仅当对任意 $x_0\in\mathring I$,存在 $\alpha\in\mathbb{R}$,使得
    \[
        f(x)\geq\alpha(x-x_0)+f(x_0).
    \]
\end{theorem}

\begin{remark}
    当 $f$ 可导时,显然可取 $\alpha=f'(x_0)$.
\end{remark}

\begin{theorem}[\markstar]
    \begin{enumerate}
        \item 一般情况下,定义 4.4.2 与定义 4.4.3 等价;
        \item $f$ 连续时,定义 4.4.1,定义 4.4.2 与定义 4.4.3 等价;
        \item $f$ 可导时,定义 4.4.1,定义 4.4.2,定义 4.4.3,定义 4.4.4 与 $f'$ 单增等价;
        \item $f$ 二阶可导时,定义 4.4.1,定义 4.4.2,定义 4.4.3,定义 4.4.4,$f'$ 单增与 $f''>0$ 等价.
    \end{enumerate}
\end{theorem}

\begin{theorem}[\markstar]
    设 $f(x)$ 为 $\lbrack a,b\rbrack$ 上的凸函数,则:
    \begin{enumerate}
        \item $f$ 在 $(a,b)$ 上连续;
        \item $f(x)\leq\max\{f(a),f(b)\}$.
    \end{enumerate}
\end{theorem}

\begin{remark}
    事实上,凸函数在定义域内部有良好性质,但边界上无法约束(就算端点特别大或特别小,亦不影响其凹凸性).
\end{remark}

%例题4.4.1
\begin{example}
    设 $f(x)$ 为区间 $(a,b)$ 内的凸函数,试证:$f(x)$ 在 $I$ 的任一闭区间 $\lbrack \alpha,\beta\rbrack\subseteq(a,b)$ 上满足 Lipschitz 条件.
\end{example}
\exampleblank

\begin{solution}
    在 $(a,b)$ 中选取 $c<\alpha<\beta<d$.由凸函数割线斜率的单调性,对任意 $\alpha\leq x<y\leq\beta$ 有
    \[
        \dfrac{f(\alpha)-f(c)}{\alpha-c}
        \leq\dfrac{f(y)-f(x)}{y-x}
        \leq\dfrac{f(d)-f(\beta)}{d-\beta}.
    \]
    令
    \[
        L=\max\left\{
        \left|\dfrac{f(\alpha)-f(c)}{\alpha-c}\right|,
        \left|\dfrac{f(d)-f(\beta)}{d-\beta}\right|
        \right\},
    \]
    便有 $|f(y)-f(x)|\leq L|y-x|$.故 $f$ 在 $\lbrack\alpha,\beta\rbrack$ 上满足 Lipschitz 条件.
\end{solution}

%例题4.4.2
\begin{example}
    设 $f(0)=0$,$f(x)$ 在 $\lbrack 0,+\infty)$ 上为非负的严格凸函数,令
    \[
        F(x)=\dfrac{f(x)}{x},\qquad x>0,
    \]
    证明:$f(x),F(x)$ 为严格递增的.
\end{example}
\exampleblank

\begin{solution}
    对 $0<x<y$,严格凸性以及 $f(0)=0$ 给出
    \[
        f(x)
        <\dfrac{x}{y}f(y)+\left(1-\dfrac{x}{y}\right)f(0)
        =\dfrac{x}{y}f(y).
    \]
    因此
    \[
        F(x)=\dfrac{f(x)}{x}<\dfrac{f(y)}{y}=F(y),
    \]
    所以 $F$ 严格递增.又因 $f$ 非负,$F(x)\geq0$,从而
    \[
        f(x)=xF(x)<yF(y)=f(y).
    \]
    故 $f$ 也严格递增.
\end{solution}

%例题4.4.3
\begin{example}
    设 $f(x),g(x)$ 在区间 $(a,b)$ 内有定义,且对任意 $x,x_0\in(a,b)$ 有
    \[
        f(x)-f(x_0)\geq g(x_0)(x-x_0).
    \]
    试证:对任意 $x_0\in(a,b)$,$f(x)$ 在 $x_0$ 处连续,并且有左右单侧导数.
\end{example}
\exampleblank

\begin{solution}
    将题设不等式中的 $x,x_0$ 交换,得到当 $x>x_0$ 时
    \[
        g(x_0)
        \leq\dfrac{f(x)-f(x_0)}{x-x_0}
        \leq g(x).
    \]
    特别地,$g$ 单调不减.在 $x_0$ 附近取 $u<x_0<v$,则该邻域内的所有差商都夹在有限数 $g(u)$ 与 $g(v)$ 之间,所以 $f$ 在 $x_0$ 处局部 Lipschitz,因而连续.

    另一方面,凸函数的右差商随 $x\downarrow x_0$ 单调,左差商随 $x\uparrow x_0$ 单调,且均被上式局部控制,故两个有限极限都存在.它们分别就是 $f'_+(x_0)$ 与 $f'_-(x_0)$.
\end{solution}

%例题4.4.4
\begin{example}
    证明下列命题成立:
    \begin{enumerate}
        \item 两凸函数之和仍为凸函数;
        \item 两递增非负凸函数之积仍为凸函数.
    \end{enumerate}
\end{example}
\exampleblank

\begin{solution}
    \begin{enumerate}
        \item 对 $0\leq\lambda\leq1$,分别对两个凸函数使用定义并相加,立即得到其和仍满足凸性不等式.

        \item 设 $x<y$,并记
              \[
                  a=f(x),\quad b=f(y),\quad c=g(x),\quad d=g(y).
              \]
              递增且非负给出 $0\leq a\leq b$,$0\leq c\leq d$.由凸性,
              \[
                  \begin{aligned}
                      f(\lambda x+(1-\lambda)y)
                      g(\lambda x+(1-\lambda)y)
                       & \leq[\lambda a+(1-\lambda)b]
                      [\lambda c+(1-\lambda)d]         \\
                       & \leq\lambda ac+(1-\lambda)bd,
                  \end{aligned}
              \]
              其中第二个不等式等价于
              \[
                  \lambda(1-\lambda)(b-a)(d-c)\geq0.
              \]
              故乘积仍为凸函数.
    \end{enumerate}
\end{solution}

%例题4.4.5
\begin{example}
    设 $f\in C(\lbrack 0,1\rbrack)$,且在 $(0,1)$ 上二次可导,若有
    \[
        f(0)=f(1)=0,
        \qquad f''(x)+2f'(x)+f(x)=0\quad(0<x<1),
    \]
    则 $f(x)\leq0$.
\end{example}
\exampleblank

\begin{solution}
    令 $u(x)=e^x f(x)$,则
    \[
        u''(x)=e^x[f''(x)+2f'(x)+f(x)]=0.
    \]
    故 $u$ 是一次函数.又 $u(0)=f(0)=0$,$u(1)=ef(1)=0$,所以 $u\equiv0$,进而 $f\equiv0$.特别地,$f(x)\leq0$.
\end{solution}

%例题4.4.6
\begin{example}
    设 $f(x)$ 在 $(a,b)$ 上四次可导,$x_0\in(a,b)$,且有
    \[
        f^{(2)}(x_0)=0,
        \qquad f^{(3)}(x_0)=0,
        \qquad f^{(4)}(x_0)>0,
    \]
    则 $f(x)$ 是下凸函数.
\end{example}
\exampleblank

\begin{solution}
    本题由一点处的导数条件不能推出整个区间上的凸性.取
    \[
        f(x)=x^4-x^6,
        \qquad x_0=0,
        \qquad (a,b)=(-2,2).
    \]
    则
    \[
        f''(0)=f'''(0)=0,
        \qquad f^{(4)}(0)=24>0,
    \]
    但
    \[
        f''(x)=12x^2-30x^4
    \]
    在区间中既取正值也取负值,所以 $f$ 在整个 $(a,b)$ 上并非凸函数.现有条件至多能在增加适当连续性后推出 $x_0$ 的某个邻域内下凸,原题结论范围需要核对.
\end{solution}

%例题4.4.7
\begin{example}
    设 $f(x)$ 在 $(-\infty,+\infty)$ 上二次可导,且有
    \[
        f(x)\leq0,
        \qquad f''(x)\geq0,
    \]
    则 $f(x)\equiv C$(常数).
\end{example}
\exampleblank

\begin{solution}
    由 $f''\geq0$,函数 $f$ 为凸函数.若存在 $x_0$ 使 $f'(x_0)>0$,则凸函数的支撑线不等式给出
    \[
        f(x)\geq f(x_0)+f'(x_0)(x-x_0),
    \]
    当 $x\to+\infty$ 时右端趋于正无穷,与 $f(x)\leq0$ 矛盾.若 $f'(x_0)<0$,令 $x\to-\infty$ 同样矛盾.因此处处有 $f'(x)=0$,故 $f$ 为常数.
\end{solution}

%例题4.4.8
\begin{example}
    设 $f(x)$ 在 $(a,b)$ 上可导,若对 $(a,b)$ 中的 $x,y$($x\neq y$),存在唯一的 $\xi\in(a,b)$,使得
    \[
        \dfrac{f(y)-f(x)}{y-x}=f'(\xi),
    \]
    则 $f(x)$ 是严格上凸或下凸的.
\end{example}
\exampleblank

\begin{solution}
    使用 Lagrange 中值定理的一个标准逆命题:若每一条割线都恰有一个中值点,则导函数必为一一映射.否则,若 $f'$ 在两点取相同值,利用导函数的 Darboux 性质并在两点之间对相应割线斜率作连续移动,可构造一条具有至少两个中值点的割线,与唯一性矛盾.

    因此 $f'$ 在区间上单射.导函数具有 Darboux 性质,所以单射的 $f'$ 必严格单调.若 $f'$ 严格递增,则 $f$ 严格下凸;若 $f'$ 严格递减,则 $f$ 严格上凸.故结论成立.
\end{solution}

%例题4.4.9
\begin{example}
    设 $f(x)$ 是 $(0,1)$ 上三次可导的非负函数,若存在 $x_1,x_2\in(0,1)$ 使得 $f(x_1)=f(x_2)=0$,则存在点 $\xi\in(0,1)$,使得
    \[
        f'''(\xi)=0.
    \]
\end{example}
\exampleblank

\begin{solution}
    不妨设 $x_1<x_2$.由于 $f\geq0$ 且 $f(x_1)=f(x_2)=0$,两点都是极小值点,故
    \[
        f'(x_1)=f'(x_2)=0.
    \]
    若 $f$ 在 $\lbrack x_1,x_2\rbrack$ 上恒为零,结论显然.否则,$f$ 在某个 $x_0\in(x_1,x_2)$ 取得正的最大值,且 $f'(x_0)=0$.于是 $f'$ 在 $x_1,x_0,x_2$ 三点为零.连续使用 Rolle 定理两次,先得到两个不同的 $f''$ 零点,再得到某个 $\xi\in(x_1,x_2)$ 使 $f'''(\xi)=0$.
\end{solution}



\subsection{不等式}

基础配套教材已经做了基本方法总结,接下来以实例来讲解具体如何运用.



\methodsection{(A) 利用单调性}

\begin{proposition}[\marktwostar]
    当 $x\in\left(0,\dfrac{\pi}{2}\right)$ 时,
    \[
        \dfrac{2}{\pi}x<\sin x<x.
    \]
\end{proposition}

%例题4.4.10
\begin{example}
    对任意 $n=1,2,3,\ldots$,求能使不等式
    \[
        \left(1+\dfrac{1}{n}\right)^{n+\alpha}
        \leq e
        \leq\left(1+\dfrac{1}{n}\right)^{n+\beta}
    \]
    成立的 $\alpha$ 的最大值与 $\beta$ 的最小值.
\end{example}
\exampleblank

\begin{solution}
    取对数后,题设不等式等价于
    \[
        \alpha\leq a_n\leq\beta,
        \qquad
        a_n=\dfrac{1}{\ln(1+\frac{1}{n})}-n.
    \]
    将 $n$ 暂视为正实数 $x$,并令
    \[
        A(x)=\dfrac{1}{\ln(1+\frac{1}{x})}-x.
    \]
    求导得
    \[
        A'(x)=\dfrac{1}{x(x+1)\left[\ln(1+\frac{1}{x})\right]^2}-1.
    \]
    对 $t>0$,函数
    \[
        \dfrac{t}{\sqrt{1+t}}-\ln(1+t)
    \]
    在 $t=0$ 处为零,且其导数为
    \[
        \dfrac{(\sqrt{1+t}-1)^2}{2(1+t)^{3/2}}>0.
    \]
    因此
    \[
        \ln\left(1+\dfrac{1}{x}\right)
        <\dfrac{1}{\sqrt{x(x+1)}},
    \]
    从而 $A'(x)>0$,即 $a_n$ 严格递增.又由
    \[
        \ln(1+t)>\dfrac{2t}{2+t}
        \qquad(t>0)
    \]
    可得 $a_n<\dfrac{1}{2}$,而 Taylor 展开给出 $a_n\to\dfrac{1}{2}$.因此
    \[
        \inf\limits_{n\geq1}a_n=a_1=\dfrac{1}{\ln2}-1,
        \qquad
        \sup\limits_{n\geq1}a_n=\dfrac{1}{2}.
    \]
    所以 $\alpha$ 的最大值与 $\beta$ 的最小值分别为
    \[
        \boxed{\alpha_{\max}=\dfrac{1}{\ln2}-1},
        \qquad
        \boxed{\beta_{\min}=\dfrac{1}{2}}.
    \]
\end{solution}

%例题4.4.11
\begin{example}
    证明如下不等式:
    \begin{enumerate}
        \item 当 $0<x_1<x_2<\dfrac{\pi}{2}$ 时,有
              \[
                  \dfrac{\tan x_2}{\tan x_1}>\dfrac{x_2}{x_1};
              \]
        \item 当 $x\in\left(0,\dfrac{\pi}{2}\right)$ 时,有
              \[
                  \dfrac{\tan x}{x}>\dfrac{x}{\sin x};
              \]
        \item 当 $x\in\left(0,\dfrac{\pi}{2}\right)$ 时,有
              \[
                  2\sin x+3\tan x>3x.
              \]
    \end{enumerate}
\end{example}
\exampleblank

\begin{solution}
    \begin{enumerate}
        \item 令 $\phi(x)=\dfrac{\tan x}{x}$.则
              \[
                  \phi'(x)=\dfrac{x\sec^2x-\tan x}{x^2}.
              \]
              分子在 $0$ 处为零,其导数为 $2x\sec^2x\tan x>0$,故 $\phi$ 严格递增,于是 $\phi(x_2)>\phi(x_1)$.

        \item 所证不等式等价于
              \[
                  \dfrac{\sin x}{x}>\sqrt{\cos x}.
              \]
              令 $\psi(x)=2\ln\left(\dfrac{\sin x}{x}\right)-\ln\cos x$.有 $\psi(0^+)=0$,且
              \[
                  \psi'(x)=2\cot x+\tan x-\dfrac{2}{x}>0,
              \]
              为验证最后一个不等式,令
              \[
                  K(x)=x(1+\cos^2x)-2\sin x\cos x.
              \]
              则 $K(0)=0$,且
              \[
                  K'(x)=\sin x\,[3\sin x-2x\cos x]>0,
              \]
              其中使用了 $\tan x>x$.故 $K(x)>0$,从而 $\psi'(x)>0$.因此 $\psi(x)>0$,结论成立.

        \item 因 $0<x<\dfrac{\pi}{2}$ 时 $\tan x>x$ 且 $\sin x>0$,所以
              \[
                  2\sin x+3\tan x>3x.
              \]
    \end{enumerate}
\end{solution}

%例题4.4.12
\begin{example}
    \begin{enumerate}
        \item 设 $b>a>e$,证明 $a^b>b^a$;
        \item 比较 $\pi^e$ 与 $e^\pi$ 的大小.
    \end{enumerate}
\end{example}
\exampleblank

\begin{solution}
    \begin{enumerate}
        \item 令 $\phi(x)=\dfrac{\ln x}{x}$.当 $x>e$ 时,
              \[
                  \phi'(x)=\dfrac{1-\ln x}{x^2}<0.
              \]
              因 $e<a<b$,故 $\dfrac{\ln a}{a}>\dfrac{\ln b}{b}$,即 $b\ln a>a\ln b$.取指数便得 $a^b>b^a$.

        \item 因 $\pi>e$,由同一单调性,
              \[
                  \dfrac{\ln\pi}{\pi}<\dfrac{\ln e}{e}=\dfrac{1}{e}.
              \]
              所以 $e\ln\pi<\pi$,从而
              \[
                  \boxed{\pi^e<e^\pi}.
              \]
    \end{enumerate}
\end{solution}



\methodsection{(B) 利用中值定理}

%例题4.4.13
\begin{example}
    证明:
    \begin{enumerate}
        \item 当 $x>0$ 时,
              \[
                  \ln\left(1+\dfrac{1}{x}\right)>\dfrac{1}{x+1};
              \]
        \item 当 $x>0$ 时,
              \[
                  \dfrac{x}{1+x^2}<\arctan x<x.
              \]
    \end{enumerate}
\end{example}
\exampleblank

\begin{solution}
    \begin{enumerate}
        \item 令
              \[
                  \phi(x)=\ln\left(1+\dfrac{1}{x}\right)-\dfrac{1}{x+1}.
              \]
              则
              \[
                  \phi'(x)=-\dfrac{1}{x(x+1)^2}<0,
                  \qquad
                  \lim\limits_{x\to+\infty}\phi(x)=0.
              \]
              因此 $\phi(x)>0$.

        \item 令 $u(x)=x-\arctan x$,则 $u(0)=0$ 且
              \[
                  u'(x)=\dfrac{x^2}{1+x^2}>0,
              \]
              故 $\arctan x<x$.再令
              \[
                  v(x)=\arctan x-\dfrac{x}{1+x^2}.
              \]
              则 $v(0)=0$ 且
              \[
                  v'(x)=\dfrac{2x^2}{(1+x^2)^2}>0.
              \]
              故 $\dfrac{x}{1+x^2}<\arctan x$.
    \end{enumerate}
\end{solution}

%例题4.4.14
\begin{example}
    \begin{enumerate}
        \item 证明:当 $s>0$ 时,
              \[
                  \dfrac{n^{s+1}}{s+1}
                  <1^s+2^s+\cdots+n^s
                  <\dfrac{(n+1)^{s+1}}{s+1};
              \]
        \item 设 $\displaystyle\lambda=\sum\limits_{k=1}^n\dfrac{1}{k}$,证明
              \[
                  e^\lambda>n+1.
              \]
    \end{enumerate}
\end{example}
\exampleblank

\begin{solution}
    \begin{enumerate}
        \item 因 $x^s$ 在 $\lbrack0,+\infty)$ 上严格递增,
              \[
                  \int_0^n x^s\,\mathrm{d}x
                  <\sum\limits_{k=1}^{n}k^s
                  <\int_1^{n+1}x^s\,\mathrm{d}x
                  <\int_0^{n+1}x^s\,\mathrm{d}x.
              \]
              计算两端积分即得题设不等式.

        \item 由 $e^t>1+t$($t>0$),
              \[
                  e^\lambda
                  =\prod\limits_{k=1}^{n}e^{\frac{1}{k}}
                  >\prod\limits_{k=1}^{n}\left(1+\dfrac{1}{k}\right)
                  =n+1.
              \]
    \end{enumerate}
\end{solution}



\methodsection{(C) 利用凹凸性}

%例题4.4.15
\begin{example}
    证明下列不等式:
    \begin{enumerate}
        \item 当 $x\in\lbrack 0,1\rbrack$ 时,
              \[
                  1+x^2\leq2^x\leq1+x;
              \]
        \item 当 $x\in\lbrack 0,1\rbrack$ 时,
              \[
                  \sin(\pi x)\leq\dfrac{\pi^2}{2}x(1-x).
              \]
    \end{enumerate}
\end{example}
\exampleblank

\begin{solution}
    \begin{enumerate}
        \item 函数 $2^x$ 为凸函数,故在 $\lbrack0,1\rbrack$ 上不超过连接端点的弦,即
              \[
                  2^x\leq1+x.
              \]
              又令 $h(x)=2^x-1-x^2$.在 $\lbrack0,1\rbrack$ 上
              \[
                  h''(x)=(\ln2)^2 2^x-2<0,
              \]
              且 $h(0)=h(1)=0$.由凹性,$h(x)\geq0$,故 $1+x^2\leq2^x$.

        \item 令
              \[
                  F(x)=\dfrac{\pi^2}{2}x(1-x)-\sin(\pi x).
              \]
              则 $F(0)=F(1)=0$,且
              \[
                  F''(x)=\pi^2[\sin(\pi x)-1]\leq0.
              \]
              因此 $F$ 为凹函数并位于端点弦的上方,即 $F(x)\geq0$.
    \end{enumerate}
\end{solution}

%例题4.4.16
\begin{example}[\markstar]
    证明下列经典不等式:
    \begin{enumerate}
        \item (平均值不等式)设 $a_i>0$,则
              \[
                  \dfrac{n}{\displaystyle\sum\limits_{i=1}^n\dfrac{1}{a_i}}
                  \leq\sqrt[n]{a_1a_2\cdots a_n}
                  \leq\dfrac{\displaystyle\sum\limits_{i=1}^n a_i}{n}.
              \]
        \item (Young 不等式)设 $a,b>0$,$0<p,q<1$,且 $\dfrac{1}{p}+\dfrac{1}{q}=1$,则
              \[
                  ab\leq\dfrac{a^p}{p}+\dfrac{b^q}{q}.
              \]
        \item (Holder 不等式)设 $a_i,b_i>0$,$0<p,q<1$,且 $\dfrac{1}{p}+\dfrac{1}{q}=1$,则
              \[
                  \sum\limits_{i=1}^n a_i b_i
                  \leq
                  \left(\sum\limits_{i=1}^n a_i^p\right)^{\frac{1}{p}}
                  \left(\sum\limits_{i=1}^n b_i^q\right)^{\frac{1}{q}}.
              \]
    \end{enumerate}
\end{example}
\exampleblank

\begin{solution}
    \begin{enumerate}
        \item 由 $\ln x$ 的凹性,
              \[
                  \dfrac{1}{n}\sum\limits_{i=1}^{n}\ln a_i
                  \leq\ln\left(\dfrac{\displaystyle\sum\limits_{i=1}^{n}a_i}{n}\right),
              \]
              取指数得到几何平均数不超过算术平均数.再对 $\dfrac{1}{a_i}$ 使用同一结论,便得调和平均数不超过几何平均数.

        \item 当前题面写成 $0<p,q<1$ 且 $\dfrac{1}{p}+\dfrac{1}{q}=1$,这样的正数 $p,q$ 不存在.Young 不等式的标准条件应为 $p,q>1$.在该条件下,由凸函数 $\dfrac{x^p}{p}$ 的切线不等式,或令 $u=a^p$,$v=b^q$ 后使用加权算术--几何平均不等式,可得
              \[
                  ab\leq\dfrac{a^p}{p}+\dfrac{b^q}{q}.
              \]

        \item 此处同样应为 $p,q>1$.令
              \[
                  A=\left(\sum\limits_{i=1}^{n}a_i^p\right)^{\frac{1}{p}},
                  \qquad
                  B=\left(\sum\limits_{i=1}^{n}b_i^q\right)^{\frac{1}{q}}.
              \]
              对 $\dfrac{a_i}{A}$ 与 $\dfrac{b_i}{B}$ 使用 Young 不等式并求和,得
              \[
                  \sum\limits_{i=1}^{n}\dfrac{a_i}{A}\dfrac{b_i}{B}
                  \leq\dfrac{1}{p}\sum\limits_{i=1}^{n}\dfrac{a_i^p}{A^p}
                  +\dfrac{1}{q}\sum\limits_{i=1}^{n}\dfrac{b_i^q}{B^q}
                  =1.
              \]
              两端乘以 $AB$ 即得 Hölder 不等式.原题两处指数范围都需要核对.
    \end{enumerate}
\end{solution}
